%% ============================================================================
%%  Connection-Consistent Planning and Volt/VAR Control of D-STATCOMs in
%%  Unbalanced Delta Feeders -- IEEE Access manuscript (ieeeaccess.cls).
%% ============================================================================
\documentclass{ieeeaccess}

\usepackage{amsmath,amssymb,amsfonts}
\usepackage{graphicx}
\usepackage{booktabs}
\usepackage{array}
\usepackage{multirow}
\usepackage{tabularx}
\newcolumntype{L}{>{\raggedright\arraybackslash}X}
\newcolumntype{C}{>{\centering\arraybackslash}X}
\newcolumntype{R}{>{\raggedleft\arraybackslash}X}
\usepackage{url}
\usepackage{xcolor}
\usepackage{textcomp}
% --- siunitx-free number/unit macros (ieeeaccess.cls clashes with siunitx) ---
\providecommand{\percent}{\%}
\newcommand{\num}[1]{#1}
\newcommand{\SI}[2]{#1\,\ensuremath{\mathrm{#2}}}
\newcommand{\numrange}[2]{#1\mbox{--}#2}
\newcommand{\SIrange}[3]{#1\mbox{--}#2\,\ensuremath{\mathrm{#3}}}
\makeatletter
% \SIlist renders a semicolon-separated list of exactly three values as
% "a, b, and c <unit>", matching siunitx list-units=single output.
\def\kie@silistiii#1;#2;#3\@nil{#1, #2, and #3}
\newcommand{\SIlist}[2]{\kie@silistiii#1\@nil\,\ensuremath{\mathrm{#2}}}
\makeatother

\usepackage{caption}
\definecolor{ieeecaptionblue}{RGB}{18,52,153}
\DeclareCaptionFont{ieeecapfont}{\fontsize{7}{8}\selectfont\sffamily\bfseries}
\DeclareCaptionFont{ieeecaplabelcolor}{\color{ieeecaptionblue}}
\captionsetup[figure]{font=ieeecapfont,labelfont=ieeecaplabelcolor,labelsep=period,justification=raggedright,singlelinecheck=false}
\captionsetup[table]{font=ieeecapfont,labelfont=ieeecaplabelcolor,labelsep=period,justification=raggedright,singlelinecheck=false}
% ieeeaccess.cls's biography environments reference a `biography' counter and a
% \if@biographyTOCentrynotmade conditional that the class never declares (IEEEtran
% uses different internal names). Declare both if missing so the no-photo
% biographies compile.
\makeatletter
\@ifundefined{c@biography}{\newcounter{biography}}{}
\expandafter\ifx\csname if@biographyTOCentrynotmade\endcsname\relax
  \newif\if@biographyTOCentrynotmade\@biographyTOCentrynotmadetrue\fi
\makeatother

\graphicspath{{figs/}}

% Convenience macros
\newcommand{\pu}{\text{ pu}}
\newcommand{\kvar}{\text{kvar}}
\newcommand{\Qmax}{Q^{\max}}
\newcommand{\Imax}{I^{\max}}
\newcommand{\Iben}{I^{\mathrm{ben}}}
\newcommand{\Vviol}{V_{\mathrm{viol}}}

% Float-placement tuning: let a float settle at the top of the column near its
% reference (usually the left column at a subsection start) instead of being
% deferred to a later right column. LaTeX two-column mode has no key that forces
% a single-column float into the left column specifically; these parameters bias
% placement toward the reference column and reduce awkward right-column stranding.
\setcounter{topnumber}{3}
\setcounter{totalnumber}{5}
\renewcommand{\topfraction}{0.92}
\renewcommand{\textfraction}{0.06}
\renewcommand{\floatpagefraction}{0.75}

\begin{document}

\history{}
\doi{}

\title{Connection-Consistent Planning and Volt/VAR Control of D-STATCOMs in
Unbalanced Delta Feeders}

\author{\uppercase{Junseok Lim}\authorrefmark{1}, \uppercase{Seungwan Kang}\authorrefmark{1}, and
\uppercase{Sungwoo Bae}\authorrefmark{1}, \IEEEmembership{Member, IEEE}}
\address[1]{Department of Electrical Engineering, Hanyang University, Seoul 04763, South Korea (e-mail: swbae@hanyang.ac.kr)}
\tfootnote{This work was supported by the Korea Institute of Energy Technology Evaluation and Planning (KETEP) and the Ministry of Climate, Energy \& Environment (MCEE) of the Republic of Korea (RS-2023-00234563, Development of power system modeling \& analysis and interoperability evaluation technology applied with grid forming based on distributed energy, and RS-2024-00422103, EV Smart Charging Platform Innovation Research Center). ORCID iDs: Junseok Lim, 0009-0004-8400-6561; Sungwoo Bae, 0000-0001-5252-1455.}

\markboth{Lim \textit{et al.}: Connection-Consistent Planning and Volt/VAR Control of D-STATCOMs}
{Lim and Bae: Connection-Consistent Planning and Volt/VAR Control of D-STATCOMs}

\corresp{Corresponding author: Sungwoo Bae (e-mail: swbae@hanyang.ac.kr).}

% ======================= ABSTRACT =======================
\begin{abstract}
On an unbalanced three-wire delta feeder, evaluating voltage on the line-to-ground
node magnitudes returned by the power-flow solver---rather than the line-to-line
voltages actually applied across the phase-to-phase load and device terminals---can
distort where a distribution static synchronous compensator (D-STATCOM) is placed
and how large it must be. This paper applies one connection-consistent voltage
evaluator---line-to-line at delta terminals and line-to-neutral at grounded-wye
terminals---uniformly across candidate screening, support selection and sizing, and
hourly dispatch. On the delta IEEE 37-node feeder the procedure returns a
verified-feasible three-device design of \SI{1041.3}{\kvar} at buses 735, 740, and
741 that raises the peak minimum terminal voltage from \num{0.8986} to
\SI{0.9565}{pu}. Under a matched enumeration-and-sizing protocol, evaluating the
same feeder on line-to-ground rather than line-to-line magnitudes increases the
returned installed capacity by \SI{23.6}{\%}, so the connection-domain choice
changes the design rather than only the reported numbers. A single-evaluator
24-hour comparison of fixed output, an equal-nominal-kvar capacitor, a
communication-free local Volt/VAR proxy, and coordinated loss-minimizing references
shows that controllable reactive dispatch reduces daily loss by
\SIrange{2.2}{5.1}{\%} relative to fixed output; the scope is steady-state,
one-quadrant reactive support, and transferability is checked on the mixed IEEE
123-node feeder.
\end{abstract}

\begin{IEEEkeywords}
D-STATCOM, unbalanced distribution feeder, delta connection, Volt/VAR control,
reactive-power planning, OpenDSS.
\end{IEEEkeywords}

\titlepgskip=-15pt
\maketitle

% ======================= NOMENCLATURE =======================
\section*{Nomenclature}
\addcontentsline{toc}{section}{Nomenclature}
\begin{IEEEdescription}[\IEEEsetlabelwidth{$P^0_{L,k},Q^0_{L,k}$~}]
\item[$i$] Index of candidate D-STATCOM buses.
\item[$m$] Index of monitored connection voltages (line-to-line for
delta buses, line-to-neutral for wye buses).
\item[$b$] Index of candidate buses in the screening stage.
\item[$t$] Hour index of the daily profile ($t=0,\dots,23$).
\item[$z_i$] Support-selection indicator, $z_i\in\{0,1\}$.
\item[$\mathcal C$] Prescreened ten-bus candidate pool.
\item[$\mathcal S$] A device support (a chosen subset of $\mathcal C$).
\item[$n_{\min},Q^\ast$] Lexicographic planning target (min count, min capacity).
\item[$\hat n,\hat{\mathbf Q}$] Smallest feasible count found and its lowest-capacity verified-feasible rating vector.
\item[$\mathcal M$] Set of monitored connection voltages.
\item[$\underline I_{i,\phi}$] Per-phase converter conductor current.
\item[$Q_i(t)$] Dispatched reactive output of the D-STATCOM at bus $i$ (kvar).
\item[$\Qmax_i$] Installed reactive rating (sizing variable) at bus $i$ (kvar).
\item[$\bar Q_i$] Per-bus reactive-capacity bound (kvar).
\item[$\underline V_{i,\phi}$] Complex phase-to-ground voltage phasor at bus $i$.
\item[$v_m$] Connection-appropriate per-unit voltage magnitude $m$.
\item[$V_{LL,i},V_{LN,i}$] Line-to-line / line-to-neutral base voltage at bus $i$ (kV).
\item[$\underline V_{i,n}$] Neutral-potential phasor at grounded-wye bus $i$.
\item[$\mathcal M_\Delta,\mathcal M_Y$] Delta and grounded-wye subsets of $\mathcal M$.
\item[$\Iben_i$] Per-device derived current benchmark at bus $i$ (from installed kvar at \SI{0.95}{pu}); a reporting reference, not a hardware limit (A).
\item[$\mathcal M_i$] Set of line-to-line terminal voltages at device bus $i$.
\item[$\Delta t$] Dispatch time step ($=\SI{1}{h}$).
\item[$\rho$] Penalty weight of the coordinated loss-minimizing reference.
\item[$\mathcal Q_{\mathrm{eval}}(S)$] Continuous rating vectors evaluated for support $S$ in the $\hat v_{\mathrm{best}}$ search.
\item[$\hat v_{\mathrm{best}}$] Largest worst-terminal minimum voltage found under the stated search.
\item[$\lambda_s,\varepsilon_s$] Near-peak stress multiplier and its Gaussian error ($\varepsilon_s\sim\mathcal N(0,0.05^2)$).
\item[$V_1,V_2$] Positive- and negative-sequence voltage magnitudes (sequence VUF).
\item[$U_{ab},U_{\mathrm{avg}}$] Line-to-line voltage magnitude and its three-phase average (NEMA LVUR).
\item[$I_\ell,\Imax_\ell$] Line current and normal ampacity of line $\ell$.
\item[$\lambda(t)$] Hourly load multiplier.
\item[$P^0_{L,k},Q^0_{L,k}$] Nominal active/reactive load at node $k$.
\item[$P_{\mathrm{loss}}$] Total feeder active-power loss (kW).
\item[$\Vviol$] Aggregate connection-voltage violation (pu).
\item[$m_v$] Planning voltage margin (pu).
\item[$\varepsilon_{\mathrm{num}}$] Numerical acceptance tolerance (pu).
\item[$E_v$] Aggregate 24-hour connection-voltage band exceedance (pu$\cdot$h).
\item[$E_{\mathrm{loss}}$] Total 24-hour energy loss (kWh).
\item[$V_i^\downarrow,V_i^\uparrow$] Lower/upper knee voltages of the local Volt/VAR proxy (pu).
\item[$\beta$] Numerical damping factor of the proxy fixed-point iteration.
\item[$\varepsilon_Q$] Convergence tolerance of the proxy iteration (kvar).
\item[$B_{\phi\psi},B_\Delta$] Per-branch and total susceptance of an equivalent delta shunt capacitor.
\item[$S_{i,\phi\psi}$] Complex power injected between phases $\phi,\psi$ at bus $i$.
\item[$\Lambda$] Screening load-multiplier set\\ $\{0.80,1.00,1.17,1.34\}$.
\item[$\Lambda_{\mathrm{plan}}$] Canonical planning snapshot\\ $\{1.34\}$ (Eq.~\eqref{eq:milp}).
\item[$R_b$] Worst-case screening voltage of bus $b$.
\end{IEEEdescription}

% ======================= INTRODUCTION =======================
\section{Introduction}
\label{sec:intro}
\IEEEPARstart{M}{aintaining} feeder voltages within an acceptable band is a binding
operational constraint under high penetration of renewable generation,
particularly on unbalanced feeders where phase asymmetry
produces uneven voltage stress \cite{keane2013dgplanning,kashem2004dgvoltage}. On a
three-wire delta feeder the loads and compensators are connected phase-to-phase, so
the voltage they actually experience is the \emph{line-to-line} magnitude, not the
line-to-ground node magnitude that a power-flow solver reports by default. When a
placement or sizing decision is made against the line-to-ground quantity, it can
misstate both the severity of the undervoltage and the reactive capacity needed to
correct it. Reactive-power support---fixed capacitors, on-load tap changers, and
distribution static synchronous compensators \cite{hingorani2000facts,singh2009statcom}---can
mitigate the voltage deterioration, but its effectiveness depends strongly on
placement, rating, and time-varying control. Throughout, this study adopts
$\numrange{0.95}{1.05}~\pu$ as an internal voltage-planning band applied to
connection-appropriate terminal voltages; it is not asserted to reproduce the
ANSI~C84.1 service-voltage categories or the IEEE~1547 point-of-common-coupling
requirements \cite{ansi_c84,ieee1547_2018}.

The device studied here is a distribution static synchronous compensator
(D-STATCOM): a voltage-source-converter reactive-support device with no
active-power exchange or energy storage in this planning model. Unlike a fixed
capacitor bank---whose injected reactive power varies with the square of the bus
voltage---the converter regulates its reactive output by control
\cite{singh2009statcom}. Consequently, \emph{installed capacity} and
\emph{dispatched output} are related but distinct decisions, and on an unbalanced
delta feeder both must be judged against the voltage the loads actually
experience---the line-to-line voltage---rather than against line-to-ground node
magnitudes, which are not the voltages applied across the delta-connected loads.
Line-to-ground quantities remain relevant to insulation and ground-fault studies,
which are outside the present voltage-compliance scope. This coupling
between the connection-appropriate voltage metric, siting, sizing, and hourly
control is the central engineering problem addressed in this paper.

Existing distribution reactive-support studies frequently treat candidate
screening, converter sizing, and time-varying control in isolation; they often
decide placement from a single voltage-stability index, model the converter
interchangeably with a capacitor, and---on delta feeders---evaluate voltages on
line-to-ground node magnitudes. This motivates the research question addressed
here: \emph{does a connection-inappropriate voltage metric materially change the
selected support location, the installed capacity, and the operating outcome, or is
it only a reporting inconsistency?} Section~\ref{sec:related} details these gaps
against recent work; the present paper answers the question within one documented
workflow on named IEEE benchmark feeders.

\subsection{Contributions}
The main contributions are as follows.
\begin{enumerate}
\item \textbf{A terminal-connection-consistent planning formulation} that applies
line-to-line voltages to phase-to-phase terminals and line-to-neutral voltages to
grounded-wye terminals throughout screening, siting, sizing, and dispatch. A
matched-metric ablation shows that line-to-ground evaluation increases the returned
installed capacity by \SI{23.6}{\%} on the IEEE 37-node feeder, changing the
engineering decision rather than only the reported numbers.
\item \textbf{An enumeration-based support-selection and two-stage sizing procedure}
that evaluates all \num{120} three-device combinations in a prescreened candidate
pool using a common full-recompilation feeder model, producing a verified-feasible
\SI{1041.3}{\kvar} design (buses 735, 740, 741) that raises the peak minimum
terminal voltage from \num{0.8986} to \SI{0.9565}{pu}.
\item \textbf{A unified 24-hour operating comparison} of fixed converter output, an
equal-nominal-kvar capacitor, a communication-free local Volt/VAR proxy, and
coordinated loss-minimizing references, quantifying \SIrange{2.2}{5.1}{\%} loss
reductions from controllable reactive-power dispatch.
\end{enumerate}

Fig.~\ref{fig:overview} summarizes the overall workflow.

\begin{figure*}[!t]
\centering
\includegraphics[width=\textwidth]{fig1_overall.pdf}
\caption{Workflow of the connection-consistent D-STATCOM framework:
constant-$Q$ system model with a line-to-line voltage metric, worst-case voltage
screening, support selection by full enumeration of the 120 three-device supports
with two-stage capacity sizing, and hourly dispatch evaluation against the
\numrange{0.95}{1.05}~pu band. The \SI{1041.3}{\kvar} vector is the lowest-capacity
verified-feasible support found under the stated enumeration budget and is carried
into the operational comparison.}
\label{fig:overview}
\end{figure*}

The remainder of the paper is organized as follows.
Section~\ref{sec:related} reviews related work and states the research gap.
Section~\ref{sec:model} formulates the system model and the placement, sizing,
and dispatch framework. Section~\ref{sec:results} reports and discusses the
simulation results, Section~\ref{sec:ieee123} tests transferability to the IEEE
123-node feeder, and Section~\ref{sec:conclusion} concludes.

% ======================= RELATED WORK =======================
\section{Related Work and Research Gap}
\label{sec:related}

\subsection{Reactive-Power Support in Distribution Systems}
Voltage-support planning in distribution systems has traditionally relied on
capacitor placement, on-load tap changers, and distribution STATCOMs
(D-STATCOMs) \cite{hingorani2000facts,edris1997facts,singh2009statcom}.
Capacitor banks are inexpensive, but their discrete, voltage-dependent operation
can be insufficient under rapidly changing load and generation
\cite{ng2000capacitor}. STATCOM-type converters instead provide fast,
voltage-independent reactive control, which is attractive for feeders with
pronounced unbalance \cite{singh2009statcom}. Reviews of D-STATCOM allocation
\cite{sirjani2017dstatcom} and comparative placement studies
\cite{taher2014dstatcom,devi2014dstatcom} confirm that the benefit of a
converter is strongly coupled to its location and rating, motivating a joint
placement-and-sizing treatment.

\subsection{Placement Criteria: Indices versus Load-Flow Evidence}
Many placement studies rank weak buses using voltage-stability indices
\cite{chakravorty2001vsi}, typically built on an efficient radial load-flow
\cite{haque1996loadflow}. Such indices compress
voltage-security information into a scalar and are convenient, but they require
careful derivation and their ranking can be sensitive to the assumed operating
point. The present work instead uses voltage screening computed directly from
unbalanced load-flow results across a range of load multipliers, treating
index-based reasoning as background rather than as the sole placement criterion,
which makes the siting justification explicit and yields a transparent,
reproducible candidate set.

\subsection{Metaheuristics for Distribution Planning}
Metaheuristic optimizers are widely used for distribution planning because the
objective is non-convex, derivative information is rarely available, and
constraints are evaluated through external simulators
\cite{keane2013dgplanning,georgilakis2013dgplacement,kennedy1995pso}. Recent
(2023--2026) allocation studies have moved toward newer and hybrid metaheuristics
and richer formulations: multi-metaheuristic STATCOM location/sizing/technology
selection \cite{statcom2025cmes}, coordinated PV and D-STATCOM compensation under
fixed versus variable reactive-power modes \cite{pvdstatcom2025fixedvar},
multi-objective Pareto integration of distributed generation and compensation
\cite{pareto2024dgcomp}, and uncertainty-aware D-STATCOM sizing
\cite{estatcom2024egwo}. Because the simulator-based continuous sizing procedure
does not provide a certified global optimality gap, the discrete three-device
supports are fully enumerated and only the ratings within each fixed support are
obtained numerically. An auxiliary five-optimizer
sensitivity check on the continuous-sizing score is reported only in the
supplementary material and neither selects nor validates the design.

\subsection{Energy-Storage-Backed Converters}
Energy-storage-backed converters and coordinated battery/PV reactive control
have been studied for voltage regulation under high DER penetration
\cite{zeraati2018bess,kabir2014pvbess}. These works
highlight the value of controllable, storage-capable converters but generally
emphasize the active-power or dynamic-control layer. The present study does not
model an energy-storage-backed converter; it isolates the steady-state
\emph{reactive}-support layer of a D-STATCOM operating in constant-$Q$ mode and
quantifies how much of the voltage-security benefit is attributable to
controllability alone. Active-power exchange and state-of-charge dynamics would
require a separate storage-enabled extension, left to future work.

On the operating side, adaptive reactive dispatch in distribution networks has
recently been pursued through receding-horizon model predictive control
\cite{rahman2024mpc,dar2024ccmpc}, multi-agent and safe deep reinforcement
learning \cite{hua2025safedrl,zheng2025marl}, and measurement-driven online
feedback optimization \cite{zhan2024ofo}. A recurring lesson from this
literature is that a time-varying dispatch should be judged against a
communication-free local baseline, because an offline schedule computed with
perfect load foresight is an upper bound rather than a deployable controller. This
paper adopts that discipline directly: in addition to the coordinated reference
and the fixed-output and capacitor baselines, Section~\ref{sec:droop} evaluates a
causal local Volt/VAR proxy (a steady-state law, without converter dynamics) so
that the reported benefit is not an artefact of perfect foresight.

\subsection{Research Gap}
The framework builds on established distribution-network and power-system modeling
\cite{baran1989reconfig,sallam2011distribution,kundur1994stability}. Based on the
studies reviewed in Section~\ref{sec:related}, three gaps remain within the scope
considered here: (i) when raw node-voltage outputs are used without
connection-aware post-processing, a three-wire delta feeder can be assessed using
line-to-ground magnitudes that are not the voltages applied across its phase-to-phase
loads; (ii) placement, sizing, and
hourly dispatch are seldom bound into one documented pipeline on an explicitly
named benchmark feeder; and (iii) the converter is often modeled interchangeably
with a capacitor, obscuring the role of controllability. The framework below
addresses these issues through connection-consistent voltage evaluation, support
selection by full enumeration with numerically constrained continuous sizing in a
prescreened pool, a device-versus-capacitor comparison, and post-design current and
thermal checks.
Table~\ref{tab:related} positions the present work against representative recent
studies on the two axes that matter here: whether the feeder is modeled as
explicitly unbalanced, and which voltage quantity the planning or control decision
is evaluated on. To avoid overclaiming, the voltage-quantity column reports
``line-to-ground'', ``line-to-line'', or ``not explicitly reported'' according to
what each source documents; it is \emph{not} asserted that a study using
line-to-ground magnitudes is in error, only that its connection-domain treatment is
not stated. The individual ingredients used here---a connection-aware metric,
support enumeration, metaheuristic sizing, and a local Volt/VAR law---are each
known; the contribution is their end-to-end integration on a named unbalanced delta
feeder and, above all, the quantitative evidence that the connection-domain choice
is not cosmetic: in the matched full-enumeration ablation, evaluating this delta
feeder on line-to-ground rather than line-to-line magnitudes raises the returned
installed capacity by \SI{23.6}{\percent} (Section~\ref{sec:certify}). The
line-to-line quantity is itself standard; what is new is enforcing it consistently
from screening through dispatch and quantifying the planning consequence of not
doing so. The source code, the OpenDSS feeder models, the input data, representative
outputs, and an automated consistency test are publicly available through the
accompanying GitHub repository (see Data Availability).

\begin{table*}[!t]
\caption{Positioning Against Representative Recent Reactive-Power Support Studies
(2024--2026)}
\label{tab:related}
\centering
\renewcommand{\arraystretch}{1.2}
\setlength{\tabcolsep}{5pt}
\begin{tabularx}{\textwidth}{@{}c L c L c c@{}}
\toprule
\textbf{Ref.} & \textbf{Feeder / connection} & \textbf{Unbalanced} &
\textbf{Voltage quantity} & \textbf{Joint siting/sizing} &
\textbf{Time-varying dispatch} \\
\midrule
\cite{statcom2025cmes}        & IEEE 14-bus system with SCIG wind & No  & Bus magnitude (not explicitly reported) & Yes & No \\
\cite{pvdstatcom2025fixedvar} & Balanced radial distribution      & No  & Bus magnitude (not explicitly reported) & Yes & Yes \\
\cite{pareto2024dgcomp}       & Balanced radial distribution      & No  & Bus magnitude (not explicitly reported) & Yes & Yes \\
\cite{estatcom2024egwo}       & Radial distribution, load unc.\   & No  & Bus magnitude (not explicitly reported) & Yes & No \\
\cite{dar2024ccmpc}           & Active distribution network       & Not explicitly reported & Bus magnitude (not explicitly reported) & No (operation) & Yes \\
\cite{zhan2024ofo}            & Distribution feeder               & Not explicitly reported & Bus magnitude (not explicitly reported) & No (operation) & Yes \\
\cite{hua2025safedrl}         & Distribution network              & Not explicitly reported & Bus magnitude (not explicitly reported) & No (operation) & Yes \\
\cite{zheng2025marl}          & Active distribution network       & Not explicitly reported & Bus magnitude (not explicitly reported) & No (operation) & Yes \\
\multirow{2}{*}{\textbf{This work}} & IEEE 37 (delta)  & Yes & Terminal-consistent L-L       & Yes & Yes \\
                                   & IEEE 123 (mixed) & Yes & Terminal-consistent L-L / L-N & No\textsuperscript{$\dagger$} & No\textsuperscript{$\dagger$} \\
\bottomrule
\multicolumn{6}{@{}p{\textwidth}}{\footnotesize ``Not explicitly reported'' marks
attributes not stated in the cited source; entries are compiled from each paper's
documented feeder model and objective and do not assert an unreported connection
domain. ``Bus magnitude'' denotes a per-unit nodal voltage evaluated without an
explicit line-to-line/line-to-neutral connection distinction.
\textsuperscript{$\dagger$}On the IEEE 123-node feeder this work performs only a
connection-domain audit and a no-action transferability check under the feeder's
native controls; joint siting/sizing and D-STATCOM dispatch are carried out on the
IEEE 37-node feeder alone.} \\
\end{tabularx}
\end{table*}

% ======================= SYSTEM MODEL =======================
\section{System Model, Problem Formulation, and Solution Framework}
\label{sec:model}
The IEEE 37-node feeder \cite{ieee37feeder} is a three-wire, delta-connected,
unbalanced \SI{4.8}{kV} distribution circuit with an open-delta head-of-feeder
voltage regulator. It is solved in OpenDSS
\cite{opendss,opendss_platform} through its Python scripting interface,
using the IEEE 37-node OpenDSS example-circuit files \cite{ieee37dss}. At each
operating point the circuit is compiled, the load multiplier is applied, the
candidate reactive sources are set, and a power-flow solution is obtained (with
the native regulator control active); the connection-appropriate per-unit
voltages defined next then form the objective terms.

\subsection{Connection-Consistent Voltage Metric}
\label{sec:connaware}
Every load, the source, and the compensator on the IEEE 37-node feeder are
connected phase-to-phase (delta); there is no neutral and no line-to-ground load.
The physically experienced voltage is therefore the \emph{line-to-line} magnitude,
and the raw line-to-ground node magnitudes returned by the solver are not the
voltages applied across the delta-connected loads, so they are not the appropriate
index for the voltage-compliance objective considered here (they remain relevant to
insulation and ground-fault studies, which are out of scope). Accordingly, each
monitored terminal $m$ at bus $i$ is evaluated on the voltage its own connection
actually applies: a line-to-line difference at delta terminals and a
line-to-neutral difference at grounded-wye terminals,
\begin{equation}
v_m =
\begin{cases}
\bigl|\underline V_{i,\phi}-\underline V_{i,\psi}\bigr|/V_{LL,i}, & m\in\mathcal M_\Delta,\\[2pt]
\bigl|\underline V_{i,\phi}-\underline V_{i,n}\bigr|/V_{LN,i}, & m\in\mathcal M_Y
\end{cases}
\label{eq:connaware}
\end{equation}
with $\phi\psi\in\{ab,bc,ca\}$ for a delta terminal, $\underline V_{i,n}$ the
neutral-potential phasor and $V_{LN,i}$ the line-to-neutral base for a grounded-wye
terminal, and $V_{LL,i}$ the line-to-line base. The monitored set is
$\mathcal M=\mathcal M_\Delta\cup\mathcal M_Y$, taken over \emph{all energized load and
device terminals}; the all-delta IEEE 37-node feeder uses only the first case, and
the same terminal-based rule supplies the line-to-neutral case for the mixed IEEE
123-node feeder (Section~\ref{sec:ieee123}). The source bus and the head-of-feeder regulator-input
bus are excluded, since no downstream shunt can raise the regulator-input voltage.
The same inclusion rule is applied in every screening, sizing, and dispatch
calculation on both feeders. On the all-delta IEEE 37-node feeder the load and
device terminals are a subset of the energized downstream buses; a scope ablation
shows that the choice between the two sets leaves the binding worst minimum
\emph{unchanged} (\SI{0.9565}{pu} at the peak, since the binding terminal is a load
bus), and changes only the light-load maximum slightly (\SI{1.032}{pu} over the
load/device terminals versus \SI{1.042}{pu} over all downstream buses), both well
inside the band. The load/device-terminal set is used throughout.
Line-to-ground node magnitudes are not the terminal voltages experienced by the
phase-to-phase loads; using \eqref{eq:connaware}
instead replaces a connection-inappropriate line-to-ground worst value of
\SI{0.88}{pu} with a line-to-line terminal minimum of \SI{0.90}{pu} (two different
voltage definitions, not the same quantity before and after correction), applied
consistently to screening, sizing, dispatch, and unbalance.

\subsection{D-STATCOM Model}
The compensator is a distribution static synchronous compensator (D-STATCOM): a
converter-interfaced controllable reactive source with no active-power exchange,
energy storage, or DC-link dynamics in this planning model. The D-STATCOM at bus
$i$ is modeled as a delta-connected constant reactive-power (Model-1) injection
with zero active power---a reactive source whose output is independent of the
local voltage within the normal operating window. In OpenDSS this is realized as a
three-phase, delta-connected generator with zero active power and a widened
constant-power voltage window of $[0.7,1.3]~\pu$
(Table~\ref{tab:settings}). This window contains every terminal voltage encountered in
screening, sizing, and dispatch---the lowest being the base-case \SI{0.8986}{pu}
line-to-line minimum---so the model does not revert to a constant-impedance
representation in any reported case. The
installed capacity $\Qmax_i$ is the converter rating reserved at bus $i$, and the
dispatched output $Q_i(t)$ is bounded by
\begin{equation}
0 \le Q_i(t) \le \Qmax_i
\label{eq:qbound}
\end{equation}
Because active power is set to zero in this planning model, $\Qmax_i$ is
interpreted as the available reactive-power rating rather than as a full
four-quadrant converter capability. Equation~\eqref{eq:qbound} restricts the
device to \emph{one-quadrant} (capacitive) reactive support; a four-quadrant
device that could also absorb reactive power ($-\Qmax_i\le Q_i\le\Qmax_i$)
subject to an apparent-power limit is left to future work. Because $\mathrm{kW}=0$,
holding a constant reactive output draws more current as the terminal voltage
falls, so a kvar bound alone does not bound the current. For a balanced
delta-connected reactive source the total output splits equally over the three
branches, $\underline S_{i,\phi\psi}(t)=jQ_i(t)/3$, and the per-phase conductor
currents follow from the branch currents,
\begin{equation}
\underline I_{i,\phi\psi}=\bigl(\underline S_{i,\phi\psi}/\underline V_{i,\phi\psi}\bigr)^{\!*},
\qquad
\underline I_{i,\phi} = \underline I_{i,\phi\psi}-\underline I_{i,\chi\phi}
\label{eq:iconv}
\end{equation}
where $\chi\phi$ denotes the branch preceding $\phi\psi$ in the cyclic order
$(a,b,c)$, so that the phase-conductor currents are
$\underline I_{i,a}=\underline I_{i,ab}-\underline I_{i,ca}$,
$\underline I_{i,b}=\underline I_{i,bc}-\underline I_{i,ab}$, and
$\underline I_{i,c}=\underline I_{i,ca}-\underline I_{i,bc}$. These
are read directly from the solved circuit; the balanced scalar
$|I_i|=|Q_i|/(\sqrt3\,|V_{LL,i}|)$ is a secondary check. For each device we define a
\emph{derived current benchmark} from its own installed output,
$\Iben_i=\Qmax_i/(\sqrt3\,V_{\!\downarrow}V_{LL,i})$ at the lower band edge
$V_{\!\downarrow}=\SI{0.95}{pu}$. Because the balanced current scales as
$|Q_i|/(\sqrt3\,V_{LL,i})$, a device sized in kvar satisfies
$\max_\phi|\underline I_{i,\phi}(t)|\le\Iben_i$ almost automatically at any voltage
above \SI{0.95}{pu}; the check is therefore a near-trivial post-design sanity
statement, \emph{not} a binding constraint and not a contribution. Consequently the
per-device currents reported in this paper are \emph{operating currents} only: no
independent hardware current or apparent-power limit is imposed in the present
planning model, and the derived benchmark is \emph{not} a nameplate-rating
verification. A genuine hardware limit would instead impose a catalog
apparent-power/current rating,
$P_i(t)^2+Q_i(t)^2\le (S_i^{\mathrm{rated}})^2$ and
$\max_\phi|\underline I_{i,\phi}(t)|\le I_i^{\mathrm{rated}}$, as an independent
\emph{input} to the sizing problem. A device whose nameplate must guarantee its
reactive rating down to the \SI{0.95}{pu} band edge would carry a nominal
apparent-power rating $S_i^{\mathrm{rated}}\ge\Qmax_i/0.95$ (for the \SI{398.5}{kvar}
bus-741 device, about \SI{419.5}{kVA}, i.e.\ \SI{50.5}{A} at the nominal
\SI{4.8}{kV}); re-sizing on such catalog ratings is the hardware-constrained
extension noted as future work. Here $P_i\equiv0$, so the reactive bound already
imposed is the only device constraint enforced.

This constant-$Q$ representation is the defining difference from a shunt
capacitor bank. A delta-connected capacitor injects, on each branch $\phi\psi$,
$Q_{\mathrm{cap},\phi\psi}(t)=B_{\phi\psi}\,|V_{\phi\psi}(t)|^2$, so its total
reactive power
\begin{equation}
Q_{\mathrm{cap},3\phi}(t) = \!\!\sum_{\phi\psi\in\{ab,bc,ca\}}\!\! B_{\phi\psi}\,|V_{\phi\psi}(t)|^2 > 0
\label{eq:cap}
\end{equation}
(with $B_{\phi\psi}=B_\Delta/3$ for a balanced bank) scales with the square of the
line-to-line voltage: it \emph{over}-injects at light load and, crucially,
\emph{under}-injects at the low-voltage peak; its OpenDSS kvar rating is the
three-phase total at rated voltage. The unbalanced
feeder is solved as a steady-state three-phase power-flow problem; for each phase
node $k$ the injected complex power relates to the nodal voltage phasor and
current through
\begin{equation}
\underline S_k(t) = P_k(t) + jQ_k(t) = \underline V_k(t)\,\underline I_k^{*}(t)
\label{eq:powerflow}
\end{equation}
where $\underline V_k$ is the complex nodal voltage; the voltage-limit and
screening terms below use the connection-appropriate magnitudes $v_m$ of
\eqref{eq:connaware} formed from these phasors. The hourly load multiplier
$\lambda(t)$ scales the nominal active and reactive loads as
\begin{equation}
P_{L,k}(t) = \lambda(t)\,P^0_{L,k}, \qquad
Q_{L,k}(t) = \lambda(t)\,Q^0_{L,k}
\label{eq:loadscale}
\end{equation}

\subsection{Objective and Placement-and-Sizing Formulation}
The voltage-limit violation aggregates the per-unit amount by which the monitored
connection voltages $v_m$ of \eqref{eq:connaware} fall below \SI{0.95}{pu} or
exceed \SI{1.05}{pu}:
\begin{equation}
\Vviol = \sum_{m}\bigl[\max(0,\,v_m-1.05) + \max(0,\,0.95-v_m)\bigr]
\label{eq:viol}
\end{equation}
The \numrange{0.95}{1.05}~pu band is applied to every monitored connection voltage
as a conservative internal planning criterion; it is distinct from service-voltage
requirements (ANSI~C84.1) and from the IEEE~1547 requirements defined at the DER
point of common coupling. The planning
\emph{target} is stated in two stages---first the smallest feasible device
count, then the least capacity at that count,
\begin{subequations}
\label{eq:milp}
\begin{align}
\text{S1:}\ & n_{\min} = \min_{\mathbf z,\mathbf Q^{\max}}\ \textstyle\sum_{i\in\mathcal C} z_i, \\
\text{S2:}\ & Q^\ast = \min_{\mathbf z,\mathbf Q^{\max}}\ \textstyle\sum_{i\in\mathcal C} \Qmax_i,\ \ \textstyle\sum_{i\in\mathcal C} z_i = n_{\min}, \\
\text{s.t.}\ & 0.95 + m_v \le v_m(\lambda) \le 1.05 - m_v, \
     \forall m\in\mathcal M,\ \lambda\in\Lambda_{\mathrm{plan}}, \\
   & 0 \le \Qmax_i \le z_i\,\bar Q_i, \quad z_i\in\{0,1\},\ \ i\in\mathcal C
\end{align}
\end{subequations}
with the candidate pool $\mathcal C$ (the prescreened ten-bus set of
Section~\ref{sec:screening}), peak-planning set $\Lambda_{\mathrm{plan}}=\{1.34\}$,
planning margin $m_v=\SI{0.006}{pu}$ (band edges \num{0.956}/\num{1.044}), and
per-bus bound $\bar Q_i=\SI{450}{\kvar}$. Equation~\eqref{eq:milp} defines the
lexicographic planning \emph{target} $(n_{\min},Q^\ast)$. The numerical procedure
returned a feasible three-device witness and found no feasible support with
$|\mathcal S|\le 2$ under the stated search protocol; we therefore report
$\hat n=3$ as the smallest feasible count found. Feasibility is checked at a
numerical tolerance $\varepsilon_{\mathrm{num}}=\ensuremath{2\times10^{-4}}~pu$ (comfortably above
the observed power-flow repeatability of $<10^{-6}$~pu), i.e.\ a design is accepted
only if $v_m\ge 0.956+\varepsilon_{\mathrm{num}}$ and
$v_m\le 1.044-\varepsilon_{\mathrm{num}}$ for all $m\in\mathcal M$ at the peak.
The selected rating vector $\hat{\mathbf Q}$ is the lowest-capacity
verified-feasible solution found under the common numerical budget and is
\emph{not} claimed to attain $Q^\ast$. The margin is
imposed only in this peak sizing problem; the 24-hour operating study
(Section~\ref{sec:operation}) reports compliance with the original
\SIrange{0.95}{1.05}{pu} band. Two quantities are deliberately \emph{outside}
\eqref{eq:milp}: the converter current \eqref{eq:iconv} is a post-design
consistency check against the derived benchmark $\Iben_i$ (obtained from $\Qmax_i$,
so it is not an independent constraint), and feeder line ampacity is a separate post-design
thermal screen (Section~\ref{sec:robust}). The discrete support choice is made by
full enumeration, but the per-support capacity is sized numerically, so
Section~\ref{sec:sizing} states the actual enumeration-plus-sizing procedure and
the scope of what it verifies.

\subsection{Two-Layer Sizing and Dispatch}
The optimization variables are interpreted in two layers. In the
\emph{planning layer}, $\Qmax_i$ is the installed converter rating at bus $i$,
sized once for the peak snapshot $\lambda_{\mathrm{peak}}=1.34$. In the
\emph{operating layer}, the per-bus dispatched outputs are chosen each hour. The
coordinated operating benchmark minimizes feeder loss subject to a
connection-voltage band $[\underline v,\overline v]$,
\begin{equation}
\mathbf{Q}(t) = \arg\!\!\min_{0\le Q_i\le \Qmax_i}\ P_{\mathrm{loss}}(t)
\ \ \text{s.t.}\ \ \underline v\le v_m(t)\le \overline v,\ \forall m
\label{eq:dispatch}
\end{equation}
approximated in implementation by the penalized objective
$P_{\mathrm{loss}}(t)+\rho\sum_m[(\underline v-v_m)_+ + (v_m-\overline v)_+]$ with
$\rho=10^{6}$, solved by bounded Powell from three starts
$\mathbf Q^{\max}\!\cdot\!\{1.0,0.6,0.3\}$ (per-coordinate tolerance $10^{-2}$,
function tolerance $10^{-3}$); the result is thus feasibility-penalized, not a
natively hard-constrained solve, and is a coordinated reference rather than a
global optimum. Because dispatch can only reduce output, the peak-planning design
is sized so the full-output case satisfies the band at $\lambda_{\mathrm{peak}}$;
feasible operation at the lighter hours is not assumed but confirmed directly by
the 24-hour evaluation of Section~\ref{sec:operation}.

\emph{Two coordinated references.}\label{sec:coordrefs}
Both operating references in this paper instantiate \eqref{eq:dispatch} with the
\emph{same} full-recompilation evaluator and the \emph{same} Powell multi-start
settings above; they differ only in the enforced band. The \emph{band-only}
reference uses the operational study band $[\underline v,\overline v]=[0.95,1.05]$~pu
introduced in Section~\ref{sec:intro}, which is an internal study band and not an
ANSI~C84.1 service-voltage category, whereas the \emph{reserve-preserving} reference
uses the tighter acceptance band $[\underline v,\overline v]=[0.9562,1.0438]$~pu that
matches the planning reserve. For both instances the final unrounded connection-voltage exceedance is
checked after optimization; the reserve-preserving reference attains a worst
connection minimum of exactly \SI{0.9562}{pu} (band exceedance verified zero),
while the band-only reference relaxes to \SI{0.9526}{pu}. Because the
communication-free proxy is compared against a constraint-matched target, the
reserve-preserving reference supplies the capture ratio of
Section~\ref{sec:operation}; the band-only reference is retained as a less
conservative operating benchmark.

% ======================= FRAMEWORK (merged into the System Model section) =======================
\label{sec:framework}
The remaining subsections turn the model above into the support-selection, sizing,
and dispatch procedure: candidate screening, support selection by enumeration, continuous
capacity sizing, exploratory optimizer comparison, post-design engineering checks,
and hourly dispatch evaluation.

\subsection{Candidate-Bus Screening}
\label{sec:screening}
Candidate buses are ranked by the worst connection voltage observed over the load
multipliers $\Lambda=\{0.80,1.00,1.17,1.34\}$, which span light-to-peak conditions
for candidate ranking only. The canonical sizing itself uses the single planning
point $\Lambda_{\mathrm{plan}}=\{1.34\}$ (Table~\ref{tab:settings}, Eq.~\eqref{eq:milp}).
The source bus and the open-delta regulator-input bus~799
are excluded, because no downstream shunt can raise the regulator-input voltage.
The screening score for bus $b$, where $M_b$ is its set of monitored connection
voltages, is
\begin{equation}
R_b = \min_{\lambda\in\Lambda}\ \min_{m\in M_b} v_m(\lambda)
\label{eq:screen}
\end{equation}
Table~\ref{tab:screening} lists the leading downstream candidates and
Fig.~\ref{fig:screening} visualizes them. The weakest buses (740, 741, 711, 738,
735, 737, and 736) form a compact low-voltage cluster with worst-case line-to-line
minima of \numrange{0.899}{0.902}~pu, below the \SI{0.95}{pu} limit; the ranking
is tight and the eighth-ranked bus~710 (\SI{0.9025}{pu}) is nearly tied, so the
number of devices actually installed is decided by the support selection of
Section~\ref{sec:sizing} rather than by a fixed cutoff. The ten weakest buses form
the candidate pool $\mathcal C$ from which the support selection draws. The pool
size ten is the low-voltage cluster itself: these are the only buses whose
worst-case minimum falls below \SI{0.91}{pu}, and the next-ranked bus is
electrically farther from the binding terminals. Enlarging the candidate pool cannot
worsen the best attainable solution and may reduce either the minimum feasible device
count or the installed capacity. Accordingly, the reported three-device result should
be interpreted as the smallest feasible count found within the prescreened ten-bus
pool under the stated search protocol; the sensitivity of the canonical capacity to
the pool size is noted as future work.

\subsection{Support Selection and Sizing Optimization}
\label{sec:sizing}
The canonical design is selected by enumerating all 120 three-device supports in
the ten-bus candidate pool. Because the pool has ten buses and the smallest
feasible count found is three (Section~\ref{sec:sensitivity}), there are only
$\binom{10}{3}=120$ three-device supports. Each fixed support is sized by a
two-stage procedure evaluated with a \emph{full feeder recompilation at every
objective evaluation}, so all 120 supports---and both voltage metrics in the
ablation of Section~\ref{sec:certify}---are compared under one identical evaluator,
the same used by the 24-hour operation. The free variables are the installed
ratings $\mathbf Q^{\max}=\{\Qmax_i\}_{i\in\mathcal S}$ of the three devices on the
fixed support $\mathcal S$ (with $\Qmax_i=0$ for $i\notin\mathcal S$), not the hourly
dispatch $Q_i(t)$. Let $[v^{\mathrm L}_s,v^{\mathrm U}_s]=[0.9567,1.0433]$~pu be the
internal search band and $[v^{\mathrm L}_a,v^{\mathrm U}_a]=[0.9562,1.0438]$~pu the
final acceptance band (a \SI{5e-4}{pu} inset on each edge). Stage~A searches for a
feasible rating vector by minimizing a quadratic penalty on \emph{both} search-band
edges,
\begin{equation}
\Phi_A(\mathbf Q^{\max})=\sum_{m\in\mathcal M}\Big[(v^{\mathrm L}_s-v_m)_+^2
+(v_m-v^{\mathrm U}_s)_+^2\Big]
\label{eq:stageA}
\end{equation}
with $v_m=v_m(\mathbf Q^{\max})$, so an over- or under-voltage start is handled
symmetrically. Stage~B, warm-started from the Stage~A point, minimizes the installed
capacity subject to the acceptance band,
\begin{equation}
\begin{aligned}
\min_{\mathbf Q^{\max}}\ &\sum_{i\in\mathcal S}\Qmax_i\\
\text{s.t.}\ &v^{\mathrm L}_a\le v_m(\mathbf Q^{\max})\le v^{\mathrm U}_a,
\quad \forall m\in\mathcal M,\\
&0\le\Qmax_i\le\SI{450}{\kvar},\ \ i\in\mathcal S,\\
&\Qmax_i=0,\ \ i\notin\mathcal S
\end{aligned}
\label{eq:stageB}
\end{equation}
Equation~\eqref{eq:stageB} states the constrained planning target; in implementation
it is \emph{numerically approximated} by the penalized objective
\begin{equation}
J_B(\mathbf Q^{\max})=\sum_{i\in\mathcal S}\Qmax_i
+\rho_V\!\sum_{m\in\mathcal M}\big[(v^{\mathrm L}_s-v_m)_+ + (v_m-v^{\mathrm U}_s)_+\big]
\label{eq:stageBimpl}
\end{equation}
with $\rho_V=10^{5}$; the penalty uses the \emph{buffered search band}, so the design
clears the acceptance band after the final \SI{0.1}{\kvar} rounding of
$\mathbf Q^{\max}$. Both stages use differential evolution (population size~8
over the three free ratings, Sobol initialization, two seeds per support; Stage~A 10
generations, Stage~B 18 generations, with local polishing). Acceptance is decided
only after that rounding and a full feeder recompilation against
$[v^{\mathrm L}_a,v^{\mathrm U}_a]$; a non-converged solve is rejected as infeasible
(none occurred in the reported runs). The lowest-capacity accepted support is $\{735,740,741\}$ at
\SI{1041.3}{\kvar}, adopted as the representative operational case by the fixed
secondary criterion of least total installed kvar among the verified-feasible
supports (ties broken by the first support returned under that rule); its nearest
rivals fall within a few kvar of this total (Section~\ref{sec:finaldesign}), so the
exact ordering among the top supports is not numerically resolved, and the
near-equivalent family is retained in Table~S3 of the supplementary material.
Enumerating the discrete supports removes any dependence on a support-search
heuristic; because each per-support capacity is sized numerically, the reported
design is the lowest-capacity verified-feasible solution found under the common
numerical budget, not a certified global optimum. As a fast cross-check, a greedy
forward selection (add the bus that most reduces the peak violation, re-sizing at
each step) recovers the same three-device count. An auxiliary five-optimizer
cross-check on the continuous-sizing score is reported only in Supplementary
Section~S-I and is not used to select the design.

\subsection{Hourly Dispatch}
The selected design is validated over a synthetic normalized 24-hour load
profile (multipliers \numrange{0.54}{1.34}) by solving the per-bus dispatch
\eqref{eq:dispatch} at each hour. The 24 operating points are solved
\emph{independently}: no intertemporal state, ramp, or switching-history constraint
is imposed, so the operational layer is an hourly steady-state dispatch evaluation
rather than an intertemporally coupled scheduling problem. The dispatch is
contrasted with two baselines:
(i) holding the same installed rating at full output every hour, and (ii)
replacing the constant-$Q$ converter with a fixed delta-connected capacitor using the
same nominal three-phase kvar parameter at \SI{1.0}{pu}. These baselines isolate, respectively, the value of
hourly controllability and the value of the constant-$Q$ device model.

\subsection{Computational Aspects}
The canonical enumeration dominates the computational cost. Each objective evaluation solves one peak-snapshot
($\lambda=1.34$) power flow; the enumeration sizes all 120 supports by the two-stage procedure
(Section~\ref{sec:sizing}) with two differential-evolution seeds each (population
size~8 over the three free ratings, i.e.\ a population of 24 vectors per
generation; Stage~A 10 generations, Stage~B 18 generations, plus local polishing),
which is on the order of $10^5$--$10^6$ power-flow solves per metric; the
line-to-line enumeration completed in about 45~minutes and the line-to-ground
enumeration in a comparable time on the reference workstation. The direct
minimum-count search and the 24-hour operating comparison add smaller amounts.
Recompiling the feeder from source at every evaluation makes the power flow a
deterministic function of the inputs, independent of evaluation order (confirmed by
re-evaluating identical inputs before and after long runs and recovering the same
voltages). The framework is implemented in Python and solves each snapshot in
OpenDSS; on a standard single-threaded workstation each full-recompile power-flow
solve takes on the order of ten milliseconds. The load-flow-based evaluation is
appropriate here precisely because placement and rating are decided before
detailed controller design, where many fast steady-state evaluations are more
valuable than a few high-fidelity dynamic simulations.

% ======================= RESULTS =======================
\section{Simulation Results and Discussion}
\label{sec:results}
The results are organized in four blocks: (A) the connection-consistent base-case
audit and candidate screening; (B) support selection and two-stage sizing; (C) the
matched line-to-ground versus line-to-line metric ablation; and (D) the 24-hour
operating comparison with its practical (unbalance, capacity-margin, and thermal)
checks. Each block is anchored on one table or figure, and quantities established in
an earlier block are not re-derived later.

\subsection{Block A: Simulation Settings and Base-Case Audit}
\label{sec:blockA}

\subsubsection{Simulation Settings}
The complete simulation settings---candidate-screening load set, the canonical
planning snapshot, the enumeration-sizing budget, and the daily operating
profile---are summarized in
Table~\ref{tab:settings}. The daily load profile is synthetic and represents
light-, mid-, and peak-load conditions; it is not measured feeder data and could
be replaced by public load data in future work. It nonetheless exposes the key
planning quantities---the peak undervoltage that sets the rating and the
off-peak hours over which controllability reduces loss.

\begin{table}[!t]
\caption{Simulation Settings}
\label{tab:settings}
\centering
\renewcommand{\arraystretch}{1.2}
\begin{tabularx}{\columnwidth}{@{}l L@{}}
\toprule
\textbf{Item} & \textbf{Setting} \\
\midrule
Feeder & IEEE 37-node unbalanced \SI{4.8}{kV} delta distribution feeder \\
Power-flow engine & OpenDSS (Python scripting interface) \\
D-STATCOM model & Delta-connected constant-$Q$ source ($\mathrm{kW}=0$) \\
Constant-power window & $[0.7,1.3]~\pu$ \\
Monitored set $\mathcal M$ & Load/device terminals, line-to-line \eqref{eq:connaware}; source, bus~799 excluded \\
Candidate pool $\mathcal C$ & 740, 741, 711, 738, 735, 737, 736, 710, 734, 733 \\
Capacity bound $\bar Q_i$ & \SI{450}{\kvar} per bus \\
Screening load set & $\lambda=\{0.80,1.00,1.17,1.34\}$ \\
Canonical planning set $\Lambda_{\mathrm{plan}}$ & $\lambda=\{1.34\}$ (peak snapshot, Eq.~\eqref{eq:milp}) \\
Canonical enum.\ sizing & 120 supports; two-stage DE (Stage A 10 gen $+$ Stage B 18 gen), population 8, Sobol init, seeds 1--2/support; full feeder recompile per evaluation; \SI{0.1}{\kvar} rounding; $\varepsilon_{\mathrm{num}}=\ensuremath{2\times10^{-4}}~pu$ \\
Min-count search & 1-D $\Delta Q=\SI{15}{\kvar}$; 2-D coarse/fine grids; final $\Delta Q=\SI{5}{\kvar}$; three Nelder--Mead cross-checks \\
Voltage limits, margin & \numrange{0.95}{1.05}~pu, $m_v=\SI{0.006}{pu}$ \\
Daily profile & Synthetic normalized 24-hour multiplier (\numrange{0.54}{1.34}) \\
Local proxy (production) & knee $(V_\downarrow,V_\uparrow)=(1.00,1.03)$~pu, $\beta=0.4$, $\varepsilon_Q=\SI{0.05}{\kvar}$, up to 400 sweeps \\
Proxy sensitivity & 27 parameter settings (9 knee $(V_\downarrow,V_\uparrow)$ pairs $\times$ 3 numerical damping $\beta$ values) \\
\bottomrule
\end{tabularx}
\end{table}

\subsubsection{Candidate-Bus Screening}
Table~\ref{tab:screening} and Fig.~\ref{fig:screening} report the worst-case
line-to-line screening voltages of \eqref{eq:connaware}. The ten candidate buses
share a tight low-voltage band (\numrange{0.899}{0.909}~pu), with bus~740 the
weakest at \SI{0.8986}{pu}. Because the ranking is numerically tight (the ten
worst-case minima span only \SI{0.010}{pu}), the number of devices is set by the support selection of
Section~\ref{sec:sizing} rather than by a fixed candidate count; the ten weakest
buses form the pool, from which the support-selection procedure of
Section~\ref{sec:sizing} selects three (735, 740, and 741). All candidate buses
lie in a contiguous downstream cluster fed through the regulator.

Two points fix the reference frame. First, the feeder-wide worst line-to-line
voltage over the day, \SI{0.8986}{pu}, occurs at bus~740, a placeable candidate;
the head-of-feeder regulator-input bus~799 is excluded because the regulator holds
its \emph{output} (bus~799r), so no downstream shunt can raise the regulator-input
voltage. Second, all magnitudes are line-to-line \eqref{eq:connaware}: with the
native regulator already at its boost limit (tap $+16$) at the peak, bus~740 still
sags to \SI{0.8986}{pu}, which is the undervoltage the D-STATCOMs address.

\begin{table}[!t]
\caption{Connection-Consistent Candidate-Bus Screening (Base Case, No Device;
Line-to-Line pu)}
\label{tab:screening}
\centering
\renewcommand{\arraystretch}{1.15}
\begin{tabularx}{\columnwidth}{@{}C C C C@{}}
\toprule
\textbf{Rank} & \textbf{Bus} & \textbf{Worst min.\ $v$ (pu)} & \textbf{Mean min.\ $v$ (pu)} \\
\midrule
1 & 740 & 0.8986 & 0.9348 \\
2 & 741 & 0.8989 & 0.9351 \\
3 & 711 & 0.8993 & 0.9354 \\
4 & 738 & 0.9005 & 0.9364 \\
5 & 735 & 0.9018 & 0.9375 \\
6 & 737 & 0.9019 & 0.9376 \\
7 & 736 & 0.9020 & 0.9377 \\
8 & 710 & 0.9025 & 0.9381 \\
9 & 734 & 0.9045 & 0.9397 \\
10 & 733 & 0.9086 & 0.9432 \\
\bottomrule
\end{tabularx}
\end{table}

\begin{figure*}[!t]
\centering
\includegraphics[width=\textwidth]{Fig2.pdf}
\caption{Worst-case minimum line-to-line voltages of the ten-bus candidate pool.
The hatched bars denote the canonical support $\{735,740,741\}$.}
\label{fig:screening}
\end{figure*}

\subsection{Block B: Support Selection and Sizing}
\label{sec:blockB}

\subsubsection{Canonical Installed Design}
\label{sec:finaldesign}
\emph{Canonical-design selection rule.} The greedy search, the multi-seed
repetition, and the metric ablation (Section~\ref{sec:certify}) are all
\emph{exploratory} sensitivity studies; the single design carried into the
operational comparison is fixed by one rule stated in advance. Because the smallest
feasible count found is three (Section~\ref{sec:sensitivity}) and only
$\binom{10}{3}=120$ three-device supports
exist, \emph{all} 120 are sized by the two-stage procedure of
Section~\ref{sec:sizing} (feasible point, then kvar minimization; seeds 1--2), with a
full feeder recompilation at every evaluation;
feasibility is judged \emph{after} the \SI{0.1}{\kvar} rounding
($v_{\min}\ge0.956+\varepsilon_{\mathrm{num}}$,
$v_{\max}\le1.044-\varepsilon_{\mathrm{num}}$, power flow converged), and
among the supports that met the
\numrange{0.9562}{1.0438}~pu acceptance the one with the lowest verified total
is taken as the representative operational case.
This yields three D-STATCOMs at buses 735, 740, and 741, listed with their derived
current benchmarks in Table~\ref{tab:design},
for a total of \SI{1041.3}{\kvar}---the lowest-capacity verified vector
found under the stated enumeration budget (population~8, two seeds per support),
narrowly ahead of its near-equivalent rivals (the next verified support $\{711,735,740\}$
sits only \SI{0.1}{\kvar} higher, and the rest within about \SI{2}{\kvar}; the exact ranking among the top
supports is not numerically resolved, Section~\ref{sec:certify}), and not a certified
global minimum. The
design raises the peak line-to-line minimum from \num{0.899} to \SI{0.9565}{pu},
i.e.\ \SI{0.0065}{pu} above the \SI{0.95}{pu} study-band lower limit, meeting the
\SI{0.006}{pu} planning reserve even after the installed capacities are rounded to
\SI{0.1}{\kvar}. The selected solution meets the \SI{0.006}{pu} planning reserve
with little additional slack because it is the lowest-capacity verified-feasible
design found under the stated numerical budget; here, \SI{0.95}{pu} is the lower
edge of the internal study band, not a service-voltage limit. It is a nominal
deterministic design at $\lambda=1.34$. Its sensitivity to load
uncertainty is characterized separately by an exploratory near-peak scalar-load stress (a Gaussian scalar-load stress of \SI{5}{\percent} standard deviation yields
band violations in \SI{7}{\percent} of realizations, Section~\ref{sec:robust}), so a
deployment margin would use scenario-based sizing or a larger reserve; this is noted
as future work. As a
reporting reference, the per-phase terminal operating currents at the peak
(\SIlist{48.2;30.6;48.8}{A}) are listed against the benchmark currents each device
would draw delivering its rating at the \SI{0.95}{pu} band edge
(\SIlist{49.8;31.6;50.5}{A}); these are operating currents, and no independent
hardware current limit is imposed in the planning model (Section~\ref{sec:model}). The direct minimum-count search
(Table~\ref{tab:sizesens}, Section~\ref{sec:sensitivity}) shows that no one- or
two-device support reached the \SI{0.9562}{pu} acceptance under the stated protocol,
so $\hat n=3$ is the smallest feasible count found under that protocol.

\begin{table}[!t]
\caption{Representative D-STATCOM Design (Verified-Feasible Three-Support Found Under the Stated Enumeration Budget):
Per-Bus Reactive Rating and Peak Operating Current}
\label{tab:design}
\centering
\renewcommand{\arraystretch}{1.2}
\setlength{\tabcolsep}{4pt}
\begin{tabularx}{\columnwidth}{@{}l *{4}{C}@{}}
\toprule
\textbf{Bus} & 735 & 740 & 741 & \textbf{Total} \\
\midrule
$\Qmax_i$ (kvar)          & 393.6 & 249.2 & 398.5 & \textbf{1041.3} \\
Peak $|I_i|$ (A)          & 48.2  & 30.6  & 48.8  & --- \\
\bottomrule
\multicolumn{5}{@{}p{0.97\columnwidth}}{\footnotesize $|I_i|$ is the peak per-phase
\emph{operating} current; no independent hardware current limit is imposed. For
reference only, the current each device would draw delivering its rating at the
\SI{0.95}{pu} band edge is $\Iben_i=\SIlist{49.8;31.6;50.5}{A}$, which the operating
currents clear---a reporting reference, not a nameplate verification.} \\
\end{tabularx}
\end{table}
The per-bus ratings are the sizing output rounded to \SI{0.1}{\kvar} (the sizing
resolution); a deployed converter would be procured at the nearest catalog rating,
and re-sizing on a discrete rating set is the natural discrete-rating extension noted
as future work. The total equals \SI{1041.3}{\kvar}.

\subsubsection{Minimum Device Count}
\label{sec:sensitivity}
The smallest feasible count is assessed by a direct numerical search rather than a
single full-output point, so it does not rest on a monotonicity assumption that a
regulator-controlled feeder need not satisfy. For every one- and two-device support
the theoretical target is the reactive dispatch that \emph{maximizes} the worst
line-to-line minimum,
\begin{equation}
v_{\mathrm{best}}(S) = \max_{0\le Q_i\le\SI{450}{\kvar}}\ \min_{m\in\mathcal M}
v_m(\mathbf Q,\lambda=1.34)
\label{eq:vbest}
\end{equation}
the search returns the largest value \emph{found} over the evaluated dispatches,
$\hat v_{\mathrm{best}}(S)=\max_{\mathbf Q\in\mathcal Q_{\mathrm{eval}}(S)}
\min_{m\in\mathcal M} v_m(\mathbf Q,\lambda=1.34)$, which lower-bounds
$v_{\mathrm{best}}(S)$. We compare $\hat v_{\mathrm{best}}$ against the planning
target \SI{0.956}{pu} (band edge \SI{0.95}{pu} plus margin $m_v$) rather than the
bare limit. One-device supports are swept on a one-dimensional grid
($\Delta Q=\SI{15}{\kvar}$, 31 points); two-device supports on a $10\times10$ coarse
grid ($\Delta Q=\SI{50}{\kvar}$) refined around the best cell to
$\Delta Q=\SI{6.2}{\kvar}$, and the single best pair is refined further to
$\Delta Q=\SI{5}{\kvar}$ and cross-checked by three independent Nelder--Mead
maximizations. Every grid point and local-search evaluation in this one- and
two-device study used the same full feeder recompilation and connection-voltage
extraction procedure as the canonical enumeration. The result (Table~\ref{tab:sizesens}) is that no $\le 2$-device
support reached the \SI{0.9562}{pu} acceptance under this grid-and-local-search
protocol: the largest value found for a single device (bus~711) was
\SI{0.9266}{pu}, and for the best pair $\{737,740\}$ only \SI{0.9490}{pu}. The best
pair's continuous cross-check returns \SI{0.9490}{pu} and its power-flow value is
repeatable to $<10^{-6}$~pu over five fresh compiles; for the one-device sweeps the
head regulators stay saturated at tap $+16$ across the whole $Q$ range. Since
$\max_S \hat v_{\mathrm{best}}(S)=\SI{0.9490}{pu}<\SI{0.956}{pu}$ for all $|S|\le2$
(a \SI{0.007}{pu} shortfall, far above the numerical resolution), no feasible
support with $|S|\le2$ was found within the candidate pool and per-bus bound under
the stated search, so three is the smallest feasible count found. The three-device canonical
design of Section~\ref{sec:finaldesign} then provides the feasibility witness
(operating minimum \SI{0.9565}{pu}; this is its dispatch value, not a
$v_{\mathrm{best}}$). Figure~\ref{fig:w4bars} summarizes the largest
worst-terminal minimum found for the one- and two-device searches and visually
distinguishes these numerical-search values from the three-device feasibility
witness. At the installed design the per-phase converter currents
are reported against the corresponding derived benchmarks (Table~\ref{tab:design});
these are operating currents, not a hardware-limit check, and the device count is set
by the voltage band.

\begin{table}[!t]
\caption{Largest Worst-Terminal Minimum Voltage $\hat v_{\mathrm{best}}$
\emph{Found Under the Stated Search} for One- and Two-Device Supports at
$\lambda=1.34$}
\label{tab:sizesens}
\centering
\renewcommand{\arraystretch}{1.2}
\setlength{\tabcolsep}{4pt}
\begin{tabularx}{\columnwidth}{@{}C L C C C@{}}
\toprule
\textbf{$n$} & \textbf{Best support} & \textbf{$\hat v_{\mathrm{best}}$ (pu)} & \textbf{$Q$ (kvar)} & \textbf{Meets \SI{0.9562}{pu}?} \\
\midrule
1 & 711      & 0.9266 & 450 & no \\
2 & 737, 740 & 0.9490 & 450, 450 & no \\
\bottomrule
\multicolumn{5}{@{}p{0.97\columnwidth}}{\footnotesize $\hat v_{\mathrm{best}}$ is the
largest value found under the grid-and-local-search protocol (a lower bound on the
true $v_{\mathrm{best}}$); $\max_{|S|\le2}\hat v_{\mathrm{best}}(S)=0.9490<0.9562$.
The separate canonical sizing yields $v_{\min}=\SI{0.9565}{pu}$ with three devices.} \\
\end{tabularx}
\end{table}

\begin{figure}[!t]
\centering
\includegraphics[width=\columnwidth]{Fig6.pdf}
\caption{Largest worst line-to-line minimum $\hat v_{\mathrm{best}}$ found under the
stated search \eqref{eq:vbest} for one- and two-device supports (bars): neither
reached the \SI{0.9562}{pu} planning target. The three-device point (diamond) is the
canonical design's \emph{operating} minimum---a feasibility witness, not a
$\hat v_{\mathrm{best}}$, so the two quantities are shown with distinct markers.}
\label{fig:w4bars}
\end{figure}


\subsection{Block C: Metric Ablation and Selection Robustness}
\label{sec:certify}
Three exploratory studies probe the design's sensitivity and isolate the effect of
the connection-consistent metric.

\emph{Selection robustness.} The canonical support is chosen by enumeration
(Section~\ref{sec:finaldesign}), not by a single stochastic run. Two facts frame the
value of the enumeration and the stability of the exact winner. First, the
enumeration is worthwhile: over the \num{120} supports for which a verified-feasible
two-stage solution was found under the stated budget, the total rating has a median of \SI{1064.2}{\kvar}, a 95th
percentile of \SI{1157.4}{\kvar}, and a maximum of \SI{1167.4}{\kvar}, so the adopted
support at \SI{1041.3}{\kvar} is about \SI{2.2}{\percent} below the median and roughly
a tenth below the worst feasible support---a saving that a feasibility-only screen
would miss. Second, the exact capacity of a support is budget-sensitive at the
\SI{0.1}{\kvar} level and the exact ranking of the top supports is \emph{not
numerically resolved}. Every support is sized with a full feeder recompilation at
each objective evaluation (Section~\ref{sec:sizing}), so the ranking is not an
artefact of a faster surrogate evaluator; under this single evaluator the
lowest-capacity accepted support is $\{735,740,741\}$ at \SI{1041.3}{\kvar}
(\SIlist{393.6;249.2;398.5}{\kvar}, $v_{\min}=\SI{0.9565}{pu}$), with the next
accepted supports (e.g.\ $\{711,735,740\}$ and $\{710,740,741\}$) within about
\SI{2}{\kvar}. The top supports are therefore numerically \emph{and practically
near-equivalent}---the few-kvar spread among them is below typical converter catalog
steps---so the reported support is retained as a representative operational case, not
because a \SI{0.1}{\kvar} numerical difference establishes engineering superiority,
and the \SI{1041.3}{\kvar} vector is the single operational source from which all
downstream results are computed.

\emph{Metric ablation.} To show that the connection-consistent metric changes the
engineering decision rather than only the reported numbers, the \emph{same} full
enumeration used for the canonical design was repeated under both metrics: all 120
three-device supports were sized under an identical per-support numerical budget
(two-stage differential evolution, seeds 1--2, full feeder recompilation per
evaluation), \SI{0.1}{\kvar} rounding, and the same
\SI{0.9562}{pu} acceptance (verified after \SI{0.1}{\kvar} rounding), once driving
feasibility on line-to-ground (LG) node magnitudes and once on line-to-line (LL)
terminal voltages (Table~\ref{tab:ablation}). The connection-inappropriate
line-to-ground diagnostic reports a \SI{0.018}{pu} lower base minimum
(\SI{0.880}{} versus \SI{0.899}{pu}), thereby \emph{overstating the apparent
undervoltage severity}. Under the matched full-enumeration and continuous-sizing
protocol, the lowest verified-feasible capacity found with the line-to-ground
objective was \SI{1287.4}{\kvar} at buses $\{711,737,740\}$
(a verified-feasible vector was found for \num{98} supports under the stated budget),
against \SI{1041.3}{\kvar}
at $\{735,740,741\}$ (\num{120} supports) with the line-to-line objective---a
\SI{23.6}{\percent} increase
($100(Q_{\mathrm{LG}}-Q_{\mathrm{LL}})/Q_{\mathrm{LL}}$). This is a
protocol-specific numerical comparison, not a global optimality certificate; because
the stochastic continuous sizing is held to a common budget, the \SI{23.6}{\percent}
difference is attributed to the evaluation metric under this matched protocol rather
than proven to isolate it exactly. As a useful check, the line-to-line enumeration
returns \emph{exactly} the canonical design. Each design satisfies its own-metric
acceptance after rounding (LG: min \SI{0.9567}{pu}, max \SI{1.0372}{pu}; LL: min
\SI{0.9565}{pu}, max \SI{1.0418}{pu}, both inside \numrange{0.9562}{1.0438}~pu),
and cross-evaluation separates over-provisioning from genuine feasibility: the
LG-metric design reaches \SI{0.9652}{pu} on the connection-consistent LL
metric---roughly \SI{0.009}{pu} of additional margin on the LL metric---while the LL
design meets the LL target at \SI{0.9565}{pu}. Evaluating the connection-inappropriate
metric therefore over-provisions the feeder by about a fifth under this protocol.

\emph{Proxy sensitivity.} The reported 24-hour local-proxy results
(Table~\ref{tab:canonical}, Figs.~\ref{fig:envelope}--\ref{fig:minv}) use the
production setting $\varepsilon_Q=\SI{0.05}{\kvar}$ with up to 400 fixed-point sweeps
\eqref{eq:proxy}; at every hour the fixed point is \emph{initialization-independent}
(starting from $Q_i^{(0)}\in\{0,\,0.5\Qmax_i,\,\Qmax_i\}$ converges to the same
dispatch within the tolerance). The knee $(V_\downarrow,V_\uparrow)$ and the damping
$\beta$ play different roles and are analyzed separately. The damping $\beta$ is used
only to stabilize the fixed-point iteration and is \emph{not} a controller parameter:
re-running $\beta\in\{0.2,0.4,0.6\}$ at the reported knee and tolerance leaves the
converged daily loss unchanged to within \SI{0.1}{\percent} of the daily total, so
the converged dispatch and loss are $\beta$-invariant in absolute terms; the
loss-\emph{capture ratio} is more sensitive only because it divides this
sub-\SI{0.1}{\percent} absolute difference by a small (\SI{170}{kWh}) denominator.
The \emph{loss capture is set by the knee}, which defines the droop curve:
sweeping $V_\downarrow\in\{0.99,1.00,1.01\}$ and $V_\uparrow\in\{1.02,1.03,1.04\}$
(all \num{27} $(V_\downarrow,V_\uparrow,\beta)$ combinations) moves the capture over
\SIrange{9.7}{79.7}{\percent} (\SI{44}{\percent} at the reported knee
$(1.00,1.03)$). All 27 settings hold the band over the whole day ($E_v=0$), so the
\emph{feasibility} conclusion is tuning-invariant while the reported capture
percentage is knee-specific.

\begin{table}[!t]
\caption{Full-Enumeration Metric Ablation: Line-to-Ground vs.\ Line-to-Line
Evaluation (Same 120-Support Enumeration Under Each Metric)}
\label{tab:ablation}
\centering
\renewcommand{\arraystretch}{1.2}
\setlength{\tabcolsep}{3pt}
\begin{tabularx}{\columnwidth}{@{}L C C C C@{}}
\toprule
\textbf{Objective} & \textbf{Support / kvar} & \textbf{Own-metric} & \textbf{Cross-metric} & \textbf{V-feas.} \\
 & & \textbf{min/max $v$} & \textbf{min/max $v$} & \textbf{found} \\
\midrule
Line-to-ground & $\{711,737,740\}$ / 1287.4 & 0.9567 / 1.0372 & 0.9652 / 1.0390 & 98 \\
Line-to-line & $\{735,740,741\}$ / 1041.3 & 0.9565 / 1.0418 & 0.9407 / 1.0355 & 120 \\
\bottomrule
\multicolumn{5}{@{}p{0.97\columnwidth}}{\footnotesize Both objectives enumerate all
120 supports under an identical per-support budget and the same
\numrange{0.9562}{1.0438}~pu acceptance (verified after \SI{0.1}{\kvar} rounding).
The last column is the number of supports for which a rounded verified-feasible
vector was \emph{found} (an algorithmic success count under the common budget, not a
certificate that the remaining supports are infeasible). Own-metric columns confirm
each design is feasible on the metric it was sized for. Relative increase
$100(Q_{\mathrm{LG}}{-}Q_{\mathrm{LL}})/Q_{\mathrm{LL}}=\SI{23.6}{\percent}$;
protocol-specific, not a global-optimality certificate.} \\
\end{tabularx}
\end{table}

\subsection{Block D: 24-Hour Operating Comparison and Practical Checks}
\label{sec:blockD}

\subsubsection{Operational Comparison of Control Strategies}
\label{sec:operation}
The installed \SI{1041.3}{\kvar} design is evaluated over the 24-hour profile by a
\emph{single} function that computes every strategy on the same model, with the
same connection-consistent violation metric \eqref{eq:viol}, and---for the coordinated
case---the same operating objective \eqref{eq:dispatch}. Five strategies are
compared: the base feeder; the design held at fixed full output; an equivalent
fixed delta-connected capacitor using the same nominal kvar parameter at \SI{1.0}{pu}; a communication-free
local voltage-droop steady-state proxy in which each converter sets its own output
from its local line-to-line terminal minimum $v_i(\mathbf Q)=\min_{m\in\mathcal M_i}
v_m(\mathbf Q)$, where $\mathcal M_i$ is the set of line-to-line terminal voltages at
device bus $i$, through the damped fixed-point iteration
\begin{equation}
Q_i^{(r+1)} = (1-\beta)\,Q_i^{(r)} + \beta\,\Qmax_i\,
\mathrm{clip}\!\left(\frac{V_i^{\uparrow}-v_i(\mathbf Q^{(r)})}{V_i^{\uparrow}-V_i^{\downarrow}},0,1\right)
\label{eq:proxy}
\end{equation}
with $V_i^{\downarrow}=1.00$ and $V_i^{\uparrow}=1.03$~pu, damping $\beta=0.4$,
initialization $Q_i^{(0)}=\Qmax_i$, the clip applied after the ratio is formed, and
the iteration stopped when $\max_i|Q_i^{(r+1)}-Q_i^{(r)}|\le\varepsilon_Q=\SI{0.05}{\kvar}$
(up to 400 sweeps); and a \emph{centrally coordinated reference}
that minimizes the feeder loss subject to the connection-voltage band
\eqref{eq:dispatch} each hour by a bounded multistart Powell search with the band
enforced by penalty (a coordinated feasibility-penalized reference, not claimed
globally optimal). The coordinated reference is evaluated in the two band forms of
Section~\ref{sec:coordrefs} (band-only and reserve-preserving), giving the six rows
of Table~\ref{tab:canonical}.
The daily violation is the connection-voltage band exceedance
$E_v=\Delta t\sum_t\sum_m[(v_m(t)-1.05)_+ + (0.95-v_m(t))_+]$ (pu$\cdot$h, $\Delta t=1$~h) and the
energy loss is $E_{\mathrm{loss}}=\Delta t\sum_t P_{\mathrm{loss}}(t)$.

Table~\ref{tab:canonical} and Figs.~\ref{fig:envelope}--\ref{fig:cap} report the
outcome; three findings stand out.

\emph{The design satisfies the study voltage band without creating overvoltage.} The constant-$Q$
design removes the base-case undervoltage: the worst line-to-line minimum rises
from \num{0.899} to \SI{0.9565}{pu} and the daily violation from $E_v=\num{4.21}$
to \num{0} (Fig.~\ref{fig:minv}). For this feeder and the selected rating, the
full-output case remained below \SI{1.05}{pu} over the evaluated light-load hours
(maximum \SI{1.032}{pu}, Fig.~\ref{fig:maxv}); this is an observed outcome of the
selected capacity and the native controls, not a general consequence of peak-based
sizing. The native regulator, already active, absorbs part of the residual excursion.

\emph{Controllability reduces losses.} With feasibility already met, the value of
hourly control is efficiency: the daily energy loss falls from \SI{2995}{kWh} at
fixed output to \SI{2928}{kWh} under droop ($-\SI{2.2}{\%}$) and \SI{2843}{kWh}
under the coordinated loss-minimizing reference ($-\SI{5.1}{\%}$), because both
inject only the reactive power each hour requires rather than full output
continuously. To keep the comparison voltage-matched, the coordinated reference is
also evaluated in a \emph{reserve-preserving} form that minimizes loss subject to the
same \numrange{0.9562}{1.0438}~pu acceptance (Section~\ref{sec:coordrefs}); its worst
minimum is \SI{0.9562}{pu} (reserve-band exceedance verified zero) versus
\SI{0.9526}{pu} for the band-only \numrange{0.95}{1.05}~pu form. The two references
are \emph{not} the same solution---their voltage and head-line-loading trajectories
differ (Table~\ref{tab:canonical})---but their unrounded daily losses agree to within
\SI{0.01}{kWh} (\SI{2843.47}{kWh} band-only versus \SI{2843.48}{kWh}
reserve-preserving). Relative to the reserve-preserving reference, and computing the
ratio from unrounded hourly losses, the local Volt/VAR proxy---using only local
measurements---recovers \SI{44.2}{\%} of the loss reduction (for the selected
$(1.00,1.03)$ knee), a substantial but less-than-half share available without
communication (Table~\ref{tab:canonical}).

\emph{The constant-$Q$ device outperforms the capacitor under the idealized device
models here.} At the same nominal three-phase kvar parameter (defined at
\SI{1.0}{pu}) the fixed capacitor \emph{satisfies the operational
\numrange{0.95}{1.05}~pu band but fails the \SI{0.006}{pu} planning reserve}
($\min v=\SI{0.9502}{pu}$, reproducible to six digits over repeated fresh
compiles---a \SI{0.0002}{pu} margin above the \SI{0.95}{pu} edge), whereas the
constant-$Q$ converter keeps \SI{0.9565}{pu}. This is because the capacitor output scales with $V^2$ (see \eqref{eq:cap}) and
\emph{under}-injects exactly when the voltage is low; it also incurs the higher
daily loss (\SI{3024}{kWh} versus \SI{2995}{kWh}, Fig.~\ref{fig:cap}) and the
higher voltage unbalance (Section~\ref{sec:vuf}). Because no independent converter
current or apparent-power limit is imposed, this is an \emph{idealized device-model}
comparison at equal nominal kvar rather than a hardware equal-nameplate test: under
that modeling assumption the voltage-independent constant-$Q$ source produces a
larger minimum-voltage margin in this feeder and profile; the comparison does not
account for a converter current limit, capacitor cost, staged switching, or step
control.

\begin{table}[!t]
\caption{Canonical 24-Hour Comparison of Control Strategies on the IEEE 37-Node
Feeder (\SI{1041.3}{\kvar} Design; Single Evaluator, Connection-Consistent Metric).
The worst minima occur at the peak hour (h19) and the maxima at light-load hours
(h4--h6); all binding terminals are line-to-line. Under the constant-$Q$ design the
worst-minimum terminal is the line-to-line pair (phases $c$--$a$) at bus~740; every
one of the 24 power flows converged (residual $<10^{-6}$~pu) and the local-proxy
iteration converged within \num{40} sweeps (of the 400 allowed) at every hour.}
\label{tab:canonical}
\centering
\renewcommand{\arraystretch}{1.2}
\setlength{\tabcolsep}{3pt}
\footnotesize
\begin{tabularx}{\columnwidth}{@{}L *{5}{C}@{}}
\toprule
\textbf{Strategy} & \textbf{Worst} & \textbf{Max.} & \textbf{$E_v$} & \textbf{Energy} & \textbf{Head-line} \\
                  & \textbf{min.\ $v$} & \textbf{$v$ (pu)} & \textbf{(pu$\cdot$h)} & \textbf{loss (kWh)} & \textbf{load.\ (\%)} \\
\midrule
Base (no device)          & 0.8986 & 1.031 & 4.213 & 3331 & 122.1 \\
Fixed max.\ output        & 0.9565 & 1.032 & 0.000 & 2995 & 114.8 \\
Capacitor (same kvar)     & 0.9502 & 1.031 & 0.000 & 3024 & 116.3 \\
Local Volt/VAR proxy      & 0.9565 & 1.031 & 0.000 & 2928 & 114.8 \\
Coordinated (band-only)   & 0.9526 & 1.032 & 0.000 & 2843 & 115.0 \\
Coordinated (reserve-pres.) & 0.9562 & 1.031 & 0.000 & 2843 & 114.8 \\
\bottomrule
\multicolumn{6}{@{}p{0.97\columnwidth}}{\footnotesize The band-only coordinated
reference enforces the operational \numrange{0.95}{1.05}~pu band \eqref{eq:dispatch}
(worst min.\ \SI{0.9526}{pu}); the reserve-preserving reference enforces the
\numrange{0.9562}{1.0438}~pu planning acceptance (worst min.\ \SI{0.9562}{pu}). The
two coordinated formulations produce different voltage and loading trajectories;
their unrounded daily losses are \SI{2843.47}{kWh} and \SI{2843.48}{kWh},
respectively, both rounding to \SI{2843}{kWh} here, so the reserve-preserving case
gives a nearly identical, voltage-matched basis for the local-proxy capture ratio.
Head-line loading is the peak-hour maximum line current as a percentage of normal
ampacity.} \\
\end{tabularx}
\end{table}

\begin{figure*}[!t]
\centering
\includegraphics[width=\textwidth]{Fig8.pdf}
\caption{24-hour line-to-line voltage envelope and dispatched reactive power under
the band-only coordinated loss-minimizing reference D-STATCOM operation. The
reserve-preserving reference, used for the constraint-matched capture ratio, is
summarized in Table~\ref{tab:canonical}.}
\label{fig:envelope}
\end{figure*}

\begin{figure*}[!t]
\centering
\includegraphics[width=\textwidth]{Fig9.pdf}
\caption{Maximum line-to-line voltage at the same nominal kvar parameter (fixed converter and
capacitor) versus the band-only coordinated reference. In this case study, all
compared strategies remain below \SI{1.05}{pu} over the evaluated light-load hours;
this is an observed result of the selected rating and feeder controls, not a general
consequence of peak-based sizing.}
\label{fig:maxv}
\end{figure*}

\begin{figure*}[!t]
\centering
\includegraphics[width=\textwidth]{Fig10.pdf}
\caption{Hourly minimum line-to-line voltage of the base feeder (worst-case
\SI{0.8986}{pu}) and of the band-only coordinated loss-minimizing reference, whose
worst minimum is \SI{0.9526}{pu} (both curves shown). The reserve-preserving
coordinated reference instead holds \SI{0.9562}{pu}, and the fixed-output and
local-proxy worst minima (\SI{0.9565}{pu}) are listed in
Table~\ref{tab:canonical}.}
\label{fig:minv}
\end{figure*}

\subsubsection{Constant-$Q$ Converter Versus Fixed Capacitor}
The capacitor and constant-$Q$ rows of Table~\ref{tab:canonical} isolate the
device model at full output. The two devices carry the \emph{same nominal
three-phase kvar parameter} (defined at \SI{1.0}{pu}); because no independent
converter current or apparent-power limit is imposed, this is an \emph{idealized
device-model} comparison rather than a hardware equal-nameplate test. The distinction
is sharper than magnitude alone suggests: the capacitor satisfies the operational
\numrange{0.95}{1.05}~pu band but fails the \SI{0.006}{pu} planning reserve
(\SI{0.9502}{pu}, only \SI{0.0002}{pu} above the \SI{0.95}{pu} edge) where the
constant-$Q$ converter keeps a comfortable
\SI{0.9565}{pu}, because the capacitor output $B_\Delta|V_{LL}|^2$ shrinks with the
square of the voltage and so \emph{under}-injects at the low-voltage peak while
\emph{over}-injecting at the high-voltage light-load hours (Fig.~\ref{fig:cap}). The
capacitor also incurs the higher daily loss (\SI{3024}{kWh} versus \SI{2995}{kWh}).
The numerical advantage therefore reflects the effect of the two steady-state device
models on this feeder and profile under the idealized constant-$Q$ assumption; it does
not account for a converter current limit, capacitor cost, staged switching, or step
control.

\begin{figure}[!t]
\centering
\includegraphics[width=\columnwidth]{Fig11.pdf}
\caption{Maximum line-to-line voltage at the same nominal kvar parameter for a fixed capacitor and a
constant-$Q$ converter over the daily profile.}
\label{fig:cap}
\end{figure}

\subsubsection{Local Volt/VAR Proxy Versus the Coordinated Reference}
\label{sec:droop}
The coordinated reference minimizes loss with full knowledge of the hourly load,
so it is an offline coordinated operating reference with perfect hourly-load
information---not a deployable controller and not a planning-layer optimization; recent
distribution voltage-control work \cite{rahman2024mpc,zhan2024ofo,hua2025safedrl}
stresses that such a schedule must be checked against a causal local law. The
local Volt/VAR law defined above serves that role: it is evaluated as a
communication-free \emph{steady-state} proxy---each converter output is the damped
fixed point (damping $\beta=0.4$, tolerance $\varepsilon_Q=\SI{0.05}{\kvar}$, up to
400 sweeps allowed; every hour converged within 40 sweeps from each tested
initialization) of its own line-to-line terminal minimum, without
measurement delay, filtering, ramp limits, or converter inner-loop dynamics, which
are outside the present steady-state scope. As Table~\ref{tab:canonical} and
Fig.~\ref{fig:droop} show, the proxy matches the coordinated reference on
feasibility---both hold $E_v=0$ over the whole day---from local measurements alone,
so the security benefit is not an artefact of perfect foresight. Coordination
retains a residual advantage in efficiency: its daily loss (\SI{2843}{kWh}) is
below the proxy's (\SI{2928}{kWh}), and, measured against the constraint-matched
reserve-preserving reference, the local law captures about \SI{44}{\%} of
the coordinated loss reduction relative to fixed output, quantifying the value of
communication. A small $(V_\downarrow,V_\uparrow,\beta)$ sensitivity
(Section~\ref{sec:certify}) confirms this share is tuning-specific but that the
feasibility conclusion is unchanged across the tested settings.

\begin{figure}[!t]
\centering
\includegraphics[width=\columnwidth]{Fig12.pdf}
\caption{Aggregate hourly dispatched reactive power $\sum_i Q_i(t)$ over the
24-hour profile for fixed full output, the local Volt/VAR proxy, and the band-only
coordinated loss-minimizing reference; the controlled strategies inject only what
each hour requires, lowering daily loss.}
\label{fig:droop}
\end{figure}

\subsubsection{Voltage-Unbalance Impact}
\label{sec:vuf}
Because the feeder is unbalanced, voltage magnitude alone does not fully capture
power quality. Table~\ref{tab:vuf} therefore adds the maximum voltage-unbalance
factor, computed from the symmetrical-component (sequence) voltages as
$\mathrm{VUF}=|V_2|/|V_1|$ from the positive- and negative-sequence voltages,
evaluated only at the energized three-phase buses (single- and two-phase buses,
where the sequence VUF is undefined, are excluded). The base feeder is strongly
unbalanced, with a worst sequence-based VUF (also termed the voltage unbalance
factor in IEEE~1159 \cite{ieee1159}) of \SI{5.27}{\percent}. Table~\ref{tab:vuf}
reports the negative-to-positive sequence ratio $|V_2|/|V_1|$ as the sequence VUF.
The NEMA~MG-1 line-voltage-unbalance rate (LVUR) is a \emph{different} metric---the
maximum line-voltage deviation from the average, divided by the average
\cite{nema_mg1}---so it is calculated and reported separately from the
line-to-line rms \emph{magnitudes} $U_{\phi\psi}=|\underline V_{\phi\psi}|$,
$\phi\psi\in\{ab,bc,ca\}$,
\begin{equation}
\begin{aligned}
U_{\mathrm{avg}}&=\tfrac13\bigl(U_{ab}+U_{bc}+U_{ca}\bigr),\\
\mathrm{LVUR}&= \frac{100\,\max_{\phi\psi}\bigl|U_{\phi\psi}-U_{\mathrm{avg}}\bigr|}
{U_{\mathrm{avg}}}
\end{aligned}
\label{eq:lvur}
\end{equation}
Because the two are not identical outside the small-unbalance regime, the
NEMA~MG-1 derating guidance is applied only to the calculated LVUR, not directly to
$|V_2|/|V_1|$, and the sequence VUF is interpreted only under the IEEE~1159
definition. NEMA~MG-1 recommends derating polyphase motors once the LVUR exceeds
about \SI{1}{\percent}; if polyphase motors were connected at the binding terminal,
the base-case LVUR of \SI{5.14}{\percent} (Table~\ref{tab:vuf}) would lie in the
range for which NEMA~MG-1 recommends derating---the present feeder model does not
explicitly simulate motor thermal derating---and the sequence VUF of
\SI{5.27}{\percent} is reported alongside it as the distinct IEEE~1159 quantity. As
Table~\ref{tab:vuf} shows, every device strategy reduces both the worst VUF (to
\SI{4.27}{\percent} at fixed output and under the local proxy, \SI{4.56}{\percent}
for the capacitor, and \SI{4.31}{\percent} under the coordinated reference) and the
worst LVUR (to \SI{4.17}{\percent}, \SI{4.42}{\percent}, and \SI{4.20}{\percent}
respectively); in this case study every device strategy was associated with a lower
maximum sequence VUF and LVUR than the base feeder, the sequence ratio falling
through network coupling rather than through direct per-phase compensation. The
slightly smaller reduction under the capacitor reflects its lower injection at the
binding hour. The residual is structural, not incidental: even at best the worst
LVUR stays near \SI{4.17}{\percent}, still far above the \SI{1}{\percent} NEMA~MG-1
guidance, because a balanced three-phase device injects equally on all phases and
cannot equalize an inherently asymmetric feeder. Removing it would require per-phase
(independently controlled) reactive support, which is outside the balanced
constant-$Q$ scope of this study.

\begin{table}[!t]
\caption{Maximum Sequence VUF and NEMA Line-Voltage-Unbalance Rate (LVUR) over the
24-Hour Profile (IEEE 37-Node, \SI{1041.3}{\kvar} Design)}
\label{tab:vuf}
\centering
\renewcommand{\arraystretch}{1.2}
\setlength{\tabcolsep}{4pt}
\begin{tabularx}{\columnwidth}{@{}l C C C@{}}
\toprule
\textbf{Strategy} & \textbf{Seq.\ VUF (\%)} & \textbf{NEMA LVUR (\%)} & \textbf{Worst min.\ $v$ (pu)} \\
\midrule
Base (no device)      & 5.27 & 5.14 & 0.8986 \\
Fixed max.\ output    & 4.27 & 4.17 & 0.9565 \\
Fixed capacitor       & 4.56 & 4.42 & 0.9502 \\
Local Volt/VAR proxy  & 4.27 & 4.17 & 0.9565 \\
Coordinated (band-only) & 4.31 & 4.20 & 0.9526 \\
\bottomrule
\multicolumn{4}{@{}p{0.94\columnwidth}}{\footnotesize Sequence VUF $=|V_2|/|V_1|$
(IEEE~1159); NEMA LVUR from \eqref{eq:lvur}. The two are distinct metrics; NEMA
derating guidance applies only to the LVUR column.} \\
\end{tabularx}
\end{table}

\subsubsection{Near-Peak Capacity-Margin Sensitivity and Thermal Screening}
\label{sec:robust}
Two post-design screens probe the planning assumptions. First, a near-peak
\emph{capacity-margin} sensitivity draws 200 realizations of a single scalar load
multiplier $\lambda_s=1.28\,(1+\varepsilon_s)$, $\varepsilon_s\sim\mathcal
N(0,0.05^2)$ (seed 2026, untruncated; all 200 power flows converged). The center
$\lambda=1.28$ is a near-peak forecast multiplier (the evening shoulder just below
the deterministic planning snapshot at $\lambda_{\mathrm{peak}}=1.34$, at which the
design gives \SI{0.9565}{pu}); this test is distinct from that planning snapshot and
is intended only to characterize margin sensitivity around the forecast. At this
stressed snapshot the local proxy reaches full output, so this experiment primarily
probes the installed \emph{capacity margin} against forecast error rather than
differences among the operating laws (which coincide when saturated). The design tolerates the error with a
residual band violation in \SI{7}{\percent} of realizations (14/200; Wilson
\SI{95}{\percent} CI \SIrange{3.7}{10.9}{\percent}); \SI{93}{\percent} of
realizations have zero violation, so the distribution is strongly right-skewed---the
95th-percentile single-snapshot $\Vviol$ (\SI{0.002}{pu}) sits just inside the
violating tail while the mean (\SI{0.003}{pu}) is lifted by a few large exceedances
(worst \SI{0.184}{pu}; these are per-snapshot band exceedances summed over the
violating terminals at that realization, e.g.\ several terminals a few hundredths
of a pu below the limit at the most stressed draw, not a daily energy).
Because $\varepsilon_s$ scales the whole feeder uniformly, this is a single-scalar
sensitivity, not a full stochastic validation---correlated, phase-specific, and
held-out scenarios are left to future work---so it bounds the peak-load capacity
margin rather than guaranteeing robustness. Second, a thermal screen confirms that voltage
feasibility does not imply thermal feasibility: at the peak ($\lambda=1.34$) the
most-loaded element is the head-of-feeder line~L35 (bus~799r to~701). It carries
\SI{488}{A} (\SI{122}{\percent} of its \SI{400}{A} normal ampacity, taken from the
feeder conductor data in the released OpenDSS model) in the base
case and \SI{459}{A} (\SI{115}{\percent}) with the D-STATCOM design held at full
output at the peak---so the reactive support reduces but does not eliminate the
pre-existing overload at this \num{1.34}$\times$ stress load. The residual head-line
loading is strategy-dependent but stays above ampacity across all operating cases
(\SIrange{114.8}{116.3}{\percent}, Table~\ref{tab:canonical}). Each installed
converter's peak per-phase terminal current (\SIlist{48.2;30.6;48.8}{A}) is
reported against its own derived benchmark (Table~\ref{tab:design}); no independent
hardware current limit is imposed. Line reinforcement is
thus a distinct, complementary planning action, treated separately from the
reactive-support design.

\subsection{Discussion}
\subsubsection{Practical Implications}
The results give two operational messages for feeder planners. First, for
phase-to-phase load-terminal voltage compliance on a delta feeder, the assessment
should use connection-appropriate line-to-line voltages: the line-to-ground node
magnitudes overstate the apparent undervoltage severity by about \SI{0.02}{pu},
reporting \SI{0.88}{pu} instead of the \SI{0.90}{pu} line-to-line terminal minimum,
and would push the sizing toward a larger design---a controlled
full-enumeration ablation (Section~\ref{sec:certify}) confirms this under the
matched enumeration-and-sizing protocol, the
connection-inappropriate line-to-ground objective returning a larger sized capacity by
about \SI{23.6}{\percent} under the matched protocol. On the corrected metric a compact \SI{1041.3}{\kvar},
three-device support---the smallest feasible count found under the stated
grid-and-local-search protocol---meets the study voltage band at
the peak (\SIrange{0.899}{0.9565}{pu}), with the per-device operating currents
reported against their derived benchmarks (no hardware current limit is modeled);
it is a \emph{voltage-feasible} reactive-support design, not a
complete network-feasible plan, because the peak head-line loading remains at
\SI{115}{\percent} of normal ampacity. Second, once the rating meets the
band, the \emph{control law} governs efficiency rather than feasibility: controlled
operation (local proxy or coordinated loss minimization) lowers the daily loss by
\SIrange{2.2}{5.1}{\percent} relative to fixed output, and the communication-free
proxy captures about \SI{44}{\%} (a tuning-specific share) of the
coordinated loss benefit. Because reconductoring
would change both impedance and voltage sensitivity, the reactive-support design
and any line reinforcement should be assessed jointly or re-optimized in sequence.

\subsubsection{Positioning Relative to the Literature}
The framework combines choices that are individually established but seldom
integrated on one named benchmark: an explicitly identified IEEE test system
solved by unbalanced load flow; a connection-consistent voltage-security metric for the
delta feeder; an enumeration-based support-selection and sizing procedure with post-design per-device
converter-current checks and separate thermal screening; and a constant-$Q$
device whose behavior is quantitatively separated from a capacitor. The placement
is returned by exhaustive support enumeration rather than by a single stochastic
solver; an auxiliary five-optimizer sensitivity check on the continuous-sizing
score (Supplementary Section~S-I) neither selects nor validates the design.

\subsubsection{Limitations and Threats to Validity}
Several limitations are stated explicitly, each pointing to a concrete extension
consistent with current practice \cite{dar2024ccmpc,pvdstatcom2025fixedvar,statcom2025cmes}.
\emph{Device model.} The study is limited to steady-state, reactive-support
(STATCOM-mode) operation; active-power dispatch, the voltage- and
active-power-dependent $P$--$Q$ capability curve, DC-link state of charge, and
converter dynamics are not modeled, so the constant-$Q$ device is a lower-fidelity
proxy for a full four-quadrant energy-storage converter. \emph{Uncertainty.} A
single deterministic 24-hour profile scaled by one scalar multiplier is used, so
it omits phase-specific, spatial, and load/PV-correlated variation; the absolute
violation magnitudes are illustrative, and a stochastic treatment with correlated,
held-out scenarios is the natural next step. The voltage dependence $Q_{\mathrm{cap}}\propto V^2$ in \eqref{eq:cap} is a general
device-model property; however, the comparative voltage margins, loss reductions,
and control benefits reported here remain specific to the feeder, load profile,
device ratings, and native control states used in this case study.
\emph{Data scope.} The profile is synthetic rather than measured; measured
load/PV data would broaden external validity. The IEEE 123-node study is a
transferability check under native controls (Section~\ref{sec:ieee123}), not an
independent second validation.
\emph{Controller tuning.} The local-proxy knee $(V_\downarrow,V_\uparrow)=(1.00,1.03)$
is a conventional Volt/VAR setting evaluated on the same profile; a held-out knee
tuning and a $(V_\downarrow,V_\uparrow,\beta)$ sensitivity are left for future work,
as are measurement delay, filtering, and converter inner-loop dynamics.
\emph{Optimality bound.} The support choice is exhaustive over the discrete space
(all 120 three-device supports), but each per-support capacity is sized by
differential evolution, so the design does not bound the distance to the global
mixed-integer optimum; a mixed-integer convex (e.g., MISOCP) relaxation would
provide such a bound.
\emph{Objective.} The installed design is fixed by the hard constraints of the
sizing problem \eqref{eq:milp}---the connection-voltage band and the per-bus kvar
bound---not by a weighted cost score; a formal multi-objective (Pareto) treatment
of the cost--security trade-off is left for future work. The auxiliary composite
score used only for the optimizer-sensitivity check is defined in Supplementary
Section~S-I.
\emph{Second feeder.} As shown next, the IEEE 123-node result is a transferability
case study conditioned on the existing regulator/capacitor controls, not an
independent generalization claim.

% ======================= IEEE 123 TRANSFERABILITY =======================
\section{Transferability to the IEEE 123-Node Feeder}
\label{sec:ieee123}
To test how the method transfers, the connection-consistent terminal audit and
voltage-security evaluation were applied to the unbalanced \SI{4.16}{kV} IEEE
123-node feeder \cite{ieee123feeder}, which has seven single-phase regulators (four
locations) and four shunt capacitors left \emph{active}
(native controls enabled), so the assessment reflects the feeder as operated. A
code audit of the released case file (Table~\ref{tab:ieee123terminals}) finds
\num{91} loads---\num{84} grounded-wye (line-to-neutral) and \num{7} delta
(line-to-line, buses 35, 65, 76)---so the connection-consistent set of
\eqref{eq:connaware} is applied \emph{per element} on this mixed feeder: the
monitored set $\mathcal M_{123}$ is the load terminals, each on its own connection
type, with the source and regulator buses 150/150r excluded.

Over the 24-hour profile, $\mathcal M_{123}$ shows \emph{no undervoltage}: the worst
connection-consistent minimum is \SI{0.9657}{pu} at wye bus~63 (line-to-neutral,
h19), well above \SI{0.95}{pu}, with the regulators at intermediate (not saturated)
taps and worst VUF only \SI{1.49}{\percent}; the worst delta terminal (bus~65) sees
a line-to-line minimum of \SI{0.9792}{pu}. This is therefore a \emph{no-action
transfer case}: because the existing regulation already holds the band, no
placement or sizing is triggered---the honest counterpart to a frozen-regulator
analysis, which would have manufactured an apparent \SI{0.898}{pu} undervoltage by
disabling the controls that resolve it. The feeder is \emph{not} fully two-sided
secure, however: a residual light-load maximum of \SI{1.052}{pu} at wye bus~83
exceeds \SI{1.05}{pu}, driven by the always-on shunt capacitors; a one-quadrant
capacitive device cannot correct it, so no complete two-sided security conclusion is
drawn. Tables~\ref{tab:ieee123} and \ref{tab:ieee123terminals} summarize the case
and the terminal-connection audit.

\begin{table*}[!t]
\begin{minipage}[t]{0.485\textwidth}
\centering
\caption{IEEE 123-Node Feeder Under Native Controls (Connection-Consistent Metric,
24-Hour Profile)}
\label{tab:ieee123}
\renewcommand{\arraystretch}{1.2}
\setlength{\tabcolsep}{4pt}
\begin{tabularx}{\linewidth}{@{}L C@{}}
\toprule
\textbf{Quantity} & \textbf{Value} \\
\midrule
Worst min.\ voltage (conn.-consistent) & \SI{0.9657}{pu} L-N (bus 63, wye, h19) \\
Worst max.\ voltage & \SI{1.052}{pu} L-N (bus 83, wye, h20) \\
Aux.\ node diagnostic (excluded) & \SI{0.9643}{pu} L-N node at delta bus 65 \\
Worst VUF           & \SI{1.49}{\percent} (bus 160, h19) \\
Regulator taps (7 units, 4 loc.) & range $-1$ to $+13$ (not saturated) \\
Shunt capacitors (4 units) & all in service \\
Additional D-STATCOM (undervoltage) & none indicated \\
\bottomrule
\end{tabularx}
\end{minipage}\hfill
\begin{minipage}[t]{0.485\textwidth}
\centering
\caption{IEEE 123-Node Terminal-Connection Audit and Per-Class Minimum (Released
Case File, Native Controls)}
\label{tab:ieee123terminals}
\renewcommand{\arraystretch}{1.2}
\setlength{\tabcolsep}{4pt}
\begin{tabularx}{\linewidth}{@{}L C C@{}}
\toprule
\textbf{Terminal class} & \textbf{Count} & \textbf{Min.\ $v$ (pu)} \\
\midrule
Grounded-wye loads (L-N) & 84 & 0.9657 (bus 63) \\
Delta loads (L-L; buses 35, 65, 76) & 7 & 0.9792 (bus 65) \\
Single-phase / three-phase loads & 89 / 2 & --- \\
Three-phase buses (VUF, seq.\ $V_2/V_1$) & 69 & --- \\
Excluded (source, 150, 150r) & 3 & --- \\
\bottomrule
\end{tabularx}
\end{minipage}
\end{table*}

Two points follow. First, the connection-consistent metric transfers across feeder
types---line-to-line on the delta IEEE 37-node feeder, line-to-neutral on the wye
IEEE 123-node feeder---selecting the physically experienced voltage in each case.
Second, because the recommendations are conditioned on the existing controls, the
output is configuration-specific: reactive support is prescribed where the controls
cannot hold the band (IEEE 37, whose regulator saturates at the peak) and declined
where they can (IEEE 123). This is a transferability and discipline check, not a
feeder-independent generalization claim; the native regulator and capacitor controls
are converged per hourly snapshot and do not reproduce chronological tap/switching
history.

% ======================= CONCLUSION =======================
\section{Conclusion}
\label{sec:conclusion}
This paper presented a connection-consistent planning and Volt/VAR-control framework
for D-STATCOMs on unbalanced delta feeders. Its central correction is that, on a
three-wire delta feeder, line-to-ground node magnitudes are not the terminal
voltages of the phase-to-phase loads, so voltage security is evaluated on
line-to-line terminal magnitudes from screening through dispatch; this replaces a
connection-inappropriate \SI{0.88}{pu} node-to-ground minimum with a \SI{0.90}{pu}
line-to-line terminal minimum. The consequence is quantitative, not cosmetic: under a
matched full-enumeration and two-stage sizing protocol, driving feasibility on
line-to-ground rather than line-to-line magnitudes raised the returned installed
capacity by \SI{23.6}{\percent}, so the connection-domain choice changes the design.
On the corrected metric, no one- or two-device support reached the \SI{0.9562}{pu}
acceptance within the ten-bus pool and \SI{450}{\kvar} per-bus bound, and enumerating
all 120 three-device supports returned a verified-feasible \SI{1041.3}{\kvar} design
at buses 735, 740, and 741 that raises the peak line-to-line minimum from \num{0.899}
to \SI{0.9565}{pu}. This is the lowest-capacity verified-feasible support found under
the stated budget, not a certified global minimum.

Once the rating meets the band, hourly controllability governs efficiency rather than
feasibility: over the deterministic 24-hour profile and a single connection-consistent
evaluator, a communication-free local proxy and a coordinated loss-minimizing
reference lowered daily loss by \SIrange{2.2}{5.1}{\percent} versus fixed output,
while an equal-nominal-kvar capacitor satisfied the operational band but failed the
\SI{0.006}{pu} planning reserve, at higher loss, because its output falls with
$V^2$. Every device case lowered the maximum voltage unbalance, though a residual near
\SI{4}{\percent} remains as a structural limit of balanced three-phase compensation.
Two checks bound the result: a near-peak scalar-load sensitivity still violated in
\SI{7}{\percent} of realizations, and the peak head-line loading stayed at
\SIrange{114.8}{116.3}{\percent} of ampacity (from \SI{122.1}{\percent} in the base
case), so the outcome is a \emph{voltage-feasible reactive-support design}, not a
complete network-feasible plan. Applied with native controls active, the same
evaluator found the regulated IEEE 123-node feeder not undervoltage-limited
(connection-consistent minimum \SI{0.9657}{pu}), a no-action transfer case. Future
work should add the converter $P$--$Q$ capability and active-power dispatch,
state-of-charge modeling, measured data, correlated uncertainty scenarios, and a
mixed-integer optimality bound.

\section*{Data Availability}
The source code, the OpenDSS feeder models, the input data, and representative output
files required to reproduce the main results---the canonical two-stage enumeration,
the metric ablation, the 24-hour operating comparison, and the IEEE 123-node
connection audit, regenerated from the listed OpenDSS models and the exact \num{24}
hourly load multipliers---are publicly available through the accompanying GitHub
repository at \url{https://github.com/Junseok-Lim-HY/dstatcom-connection-consistent}.
The version used for this article is archived under Release
\texttt{v1.0-ieee-access-submission}. The per-script roles and the artifacts checked
are itemized in Supplementary Table~S4 and the repository README. An automated
consistency test verifies cross-file numerical identities against the stored outputs
and is \emph{not} a substitute for rerunning the optimization. Only transient solver
logs and large intermediate power-flow dumps are withheld and are not required to
reproduce any reported result. The IEEE 37- and 123-node feeder OpenDSS models are the
IEEE PES distribution test feeders \cite{ieee37feeder,ieee123feeder} as distributed
with OpenDSS \cite{opendss}.

% ======================= REFERENCES =======================
\bibliographystyle{IEEEtran}
\bibliography{references}

\begin{IEEEbiography}[{\includegraphics[width=1in,height=1.25in,clip,keepaspectratio]{jslim.jpg}}]{JUNSEOK LIM}
received the B.S. degree in electrical engineering from Hanyang University,
Ansan, South Korea, in 2019. He is currently working toward the Ph.D.
degree in the Department of Electrical Engineering at Hanyang University.
His research interests include power system planning and operation,
energy storage integration, and resilience enhancement in renewable-rich
distribution networks.
\end{IEEEbiography}

\begin{IEEEbiographynophoto}{SEUNGWAN KANG}
received the B.S. and M.S. degrees in chemical engineering from Hanyang University,
Seoul, South Korea. He is currently working toward the Ph.D. degree in the Department
of Electrical Engineering at Hanyang University, Seoul, South Korea. He is also a
Director with Korea South-East Power Company, South Korea.
\end{IEEEbiographynophoto}

\begin{IEEEbiography}[{\includegraphics[width=1in,height=1.25in,clip,keepaspectratio]{swbae.jpg}}]{SUNGWOO BAE}
(Member, IEEE) received the B.S. degree from Hanyang University, Seoul,
Korea, and the M.S.E. and Ph.D. degrees from the University of Texas at
Austin, USA, all in electrical engineering, in 2006, 2009, and 2011,
respectively. From 2012 to 2013, he was a senior research engineer with
Power Center at Samsung Advanced Institute of Technology. From 2013 to
2017, he was with Yeungnam University. He has been a Professor in the
Department of Electrical Engineering at Hanyang University in Korea
since 2017. In 2005, Dr.\ Bae was awarded the Grand Prize at the
national electrical engineering design contest by the Minister of
Commerce, Industry and Energy of the Republic of Korea.
\end{IEEEbiography}

\EOD

\end{document}
